% Options for packages loaded elsewhere
\PassOptionsToPackage{unicode}{hyperref}
\PassOptionsToPackage{hyphens}{url}
\documentclass[
  11pt,
]{article}
\usepackage{xcolor}
\usepackage[left=3cm,right=5cm,top=2cm,bottom=2cm]{geometry}
\usepackage{amsmath,amssymb}
\setcounter{secnumdepth}{5}
\usepackage{iftex}
\ifPDFTeX
  \usepackage[T1]{fontenc}
  \usepackage[utf8]{inputenc}
  \usepackage{textcomp} % provide euro and other symbols
\else % if luatex or xetex
  \usepackage{unicode-math} % this also loads fontspec
  \defaultfontfeatures{Scale=MatchLowercase}
  \defaultfontfeatures[\rmfamily]{Ligatures=TeX,Scale=1}
\fi
\usepackage{lmodern}
\ifPDFTeX\else
  % xetex/luatex font selection
\fi
% Use upquote if available, for straight quotes in verbatim environments
\IfFileExists{upquote.sty}{\usepackage{upquote}}{}
\IfFileExists{microtype.sty}{% use microtype if available
  \usepackage[]{microtype}
  \UseMicrotypeSet[protrusion]{basicmath} % disable protrusion for tt fonts
}{}
\makeatletter
\@ifundefined{KOMAClassName}{% if non-KOMA class
  \IfFileExists{parskip.sty}{%
    \usepackage{parskip}
  }{% else
    \setlength{\parindent}{0pt}
    \setlength{\parskip}{6pt plus 2pt minus 1pt}}
}{% if KOMA class
  \KOMAoptions{parskip=half}}
\makeatother
\usepackage{longtable,booktabs,array}
\usepackage{caption}
\captionsetup[table]{skip=6pt}
\newcounter{none} % for unnumbered tables
\usepackage{calc} % for calculating minipage widths
% Correct order of tables after \paragraph or \subparagraph
\usepackage{etoolbox}
\makeatletter
\patchcmd\longtable{\par}{\if@noskipsec\mbox{}\fi\par}{}{}
\makeatother
% Allow footnotes in longtable head/foot
\IfFileExists{footnotehyper.sty}{\usepackage{footnotehyper}}{\usepackage{footnote}}
\makesavenoteenv{longtable}
\usepackage{graphicx}
\makeatletter
\newsavebox\pandoc@box
\newcommand*\pandocbounded[1]{% scales image to fit in text height/width
  \sbox\pandoc@box{#1}%
  \Gscale@div\@tempa{\textheight}{\dimexpr\ht\pandoc@box+\dp\pandoc@box\relax}%
  \Gscale@div\@tempb{\linewidth}{\wd\pandoc@box}%
  \ifdim\@tempb\p@<\@tempa\p@\let\@tempa\@tempb\fi% select the smaller of both
  \ifdim\@tempa\p@<\p@\scalebox{\@tempa}{\usebox\pandoc@box}%
  \else\usebox{\pandoc@box}%
  \fi%
}
% Set default figure placement to htbp
\def\fps@figure{htbp}
\makeatother
% definitions for citeproc citations
\NewDocumentCommand\citeproctext{}{}
\NewDocumentCommand\citeproc{mm}{%
  \begingroup\def\citeproctext{#2}\cite{#1}\endgroup}
\makeatletter
 % allow citations to break across lines
 \let\@cite@ofmt\@firstofone
 % avoid brackets around text for \cite:
 \def\@biblabel#1{}
 \def\@cite#1#2{{#1\if@tempswa , #2\fi}}
\makeatother
\newlength{\cslhangindent}
\setlength{\cslhangindent}{1.5em}
\newlength{\csllabelwidth}
\setlength{\csllabelwidth}{3em}
\newenvironment{CSLReferences}[2] % #1 hanging-indent, #2 entry-spacing
 {\begin{list}{}{%
  \setlength{\itemindent}{0pt}
  \setlength{\leftmargin}{0pt}
  \setlength{\parsep}{0pt}
  % turn on hanging indent if param 1 is 1
  \ifodd #1
   \setlength{\leftmargin}{\cslhangindent}
   \setlength{\itemindent}{-1\cslhangindent}
  \fi
  % set entry spacing
  \setlength{\itemsep}{#2\baselineskip}}}
 {\end{list}}
\usepackage{calc}
\newcommand{\CSLBlock}[1]{\hfill\break\parbox[t]{\linewidth}{\strut\ignorespaces#1\strut}}
\newcommand{\CSLLeftMargin}[1]{\parbox[t]{\csllabelwidth}{\strut#1\strut}}
\newcommand{\CSLRightInline}[1]{\parbox[t]{\linewidth - \csllabelwidth}{\strut#1\strut}}
\newcommand{\CSLIndent}[1]{\hspace{\cslhangindent}#1}
\setlength{\emergencystretch}{3em} % prevent overfull lines
\providecommand{\tightlist}{%
  \setlength{\itemsep}{0pt}\setlength{\parskip}{0pt}}
% Statistics in Medicine submission preamble.
% Shared across analysis/report/<paper>/report.Rmd files via
%   includes:
%     in_header: ../sim-preamble.tex
% in the YAML output block.
%
% The page footer (source path, render time, version) is no
% longer set here. It is supplied at render time by the vendored
% tools/stamp-render.R, which injects tools/stamp.tex.

% Continuous line numbering for editorial review. The 'mathlines'
% option extends numbering through display-math environments
% (equation, align, ...), which SIM editorial workflows expect.
\usepackage[mathlines]{lineno}
\linenumbers
\usepackage{amsmath}

% Spacing: linestretch is set in YAML (1.5 line spacing).
\usepackage{setspace}
% Three-part paragraph scaffolding for manuscript revision.
%
% bullets  -- boiled-down topic sentence + bullet points
%             small, sans-serif, dark gray; paragraph number [N] prepended
% rgt      -- author's rewrite (initially a placeholder)
%             small, italic, dark blue
% orig     -- original Claude-authored draft
%             normal size, light gray
%
% Presentation order per paragraph: bullets -> rgt -> orig
% Strip with: strip-claudecode -i report.Rmd

\usepackage{xcolor}

\newcounter{pgraph}

\newenvironment{bullets}{%
  \stepcounter{pgraph}%
  \begin{quote}\small\sffamily\color[gray]{0.30}%
  {\normalfont\scriptsize\color[gray]{0.58}[\thepgraph]\enspace}}{%
  \end{quote}}

\newenvironment{rgt}{%
  \begin{quote}\small\itshape\color[RGB]{20,60,160}}{%
  \end{quote}}

\newenvironment{orig}{%
  \begin{quote}\normalsize\color[gray]{0.55}}{%
  \end{quote}}

% backward compat alias
\newenvironment{claudecode}{\begin{orig}}{\end{orig}}
\usepackage{amssymb}
\usepackage{lineno}
\linenumbers
\usepackage{setspace}
\setstretch{1.1}
\usepackage{bookmark}
\IfFileExists{xurl.sty}{\usepackage{xurl}}{} % add URL line breaks if available
\urlstyle{same}
% fallback for those not using the hyperref driver hyperxmp:
\makeatletter
\@ifundefined{xmpquote}{\newcommand{\xmpquote}[1]{#1}}{}
\makeatother
\hypersetup{
  pdftitle={Latent-class and mixture-model formulations for biomarker-treatment interaction in N-of-1 trials},
  pdfauthor={pmsimstats team},
  hidelinks,
  pdfcreator={LaTeX via pandoc}}

\title{Latent-class and mixture-model formulations for
biomarker-treatment interaction in N-of-1 trials}
\author{pmsimstats team}
\date{2026-07-30}

\begin{document}
\maketitle

{
\setcounter{tocdepth}{2}
\tableofcontents
}
\section{Abstract}\label{abstract}

\textbf{Background.} Aggregated N-of-1 clinical trials test
biomarker-treatment interactions through a fixed continuous-
heterogeneity regression in published practice. The psychometric
literature has developed mature finite-mixture and latent-class
machinery for representing discrete subpopulation structure that the
biostatistical literature on N-of-1 trials has not engaged. The choice
between continuous-heterogeneity and latent-class encodings of the
biomarker-moderation question is a substantive scientific question
rather than a parameterization convenience, but the field has not
characterized when class-aware analyses become competitive with
continuous regression at trial-relevant sample sizes.

\textbf{Methods.} We develop a five-component biomarker-moderation
spectrum (covariance-only, mean-and-covariance, location-scale,
heterogeneous random-effects, latent-class) bridging
continuous-heterogeneity and latent-class formulations, and we catalogue
the finite-mixture taxonomy (latent class analysis, growth mixture
models, factor mixture models, regression mixtures) for application to
within-subject longitudinal designs. We commit a
heterogeneous-random-slopes generative data-generating process (DGP) and
a class-aware analysis pipeline (\texttt{lcmm::hlme} with gridsearch
initialization), and run a 240-replicate-per-cell prototype simulation
at \(N = 70\), gating slope \(\in \{0, 1.0, 1.5\}\), on a 4-core
parallel implementation. Three test forms are recorded per replicate:
the linear-mixed \texttt{bm:Dbc} Wald, the unconditional Wald on the
lcmm gating coefficient, and a BIC-conditional gating procedure.

\textbf{Results.} In a 240-replicate-per-cell pilot we find that the
unconditional Wald test on the latent-class gating coefficient inflates
Type I error to \(0.197\) (95\% CI \([0.145, 0.250]\)) at the null,
decisively above nominal \(0.05\) and consistent with
the\textsuperscript{\citeproc{ref-bauercurran2003}{1}}
and\textsuperscript{\citeproc{ref-lubkemuthen2005}{2}} overextraction
predictions. The BIC-conditional gating procedure controls Type I
(\(0.004\) under the null) but is severely under-powered at
trial-relevant \(N\) (rejection rate \(0.025\) at gating slope \(1.0\),
\(0.004\) at \(1.5\)). The linear-mixed \texttt{bm:Dbc} Wald test holds
nominal Type I and reaches \(0.542\) power at gating slope \(1.5\),
dominating both class-aware tests. The Monte Carlo standard error (MCSE)
on power is approximately \(0.030\) across the grid. These pilot
estimates precede the production five-study program described in the
companion simulation plan, which has not yet been executed.

\textbf{Conclusions.} Within the 240-replicate pilot, at trial- relevant
sample sizes (\(N = 70\), eight measurement timepoints), neither the
unconditional Wald-on-gating nor the BIC-conditional gating test is
competitive with the linear-mixed baseline under the
heterogeneous-random-slopes DGP. The latent-class encoding appears to
require a regime of large class separation for class-aware tests to gain
power; the pre-registered production program (five studies,
\(N \in \{35, 70, 100\}\), wider gating- slope grid) has not yet been
executed and is required to map the regime where the BIC-conditional
procedure becomes competitive.

\section{Introduction}\label{introduction}

\subsection{The aggregated N-of-1 trial and the biomarker-moderation
question}\label{the-aggregated-n-of-1-trial-and-the-biomarker-moderation-question}

\begin{bullets}

\textbf{The aggregated N-of-1 design pools within-person treatment
contrasts across participants for population-level inference.}

\begin{itemize}\setlength{\itemsep}{2pt}\setlength{\parskip}{0pt}
\item Each participant cycles between active treatment and control
  across multiple within-person blocks.
\item The design outperforms parallel-group trials in information per
  participant when treatment response is reversible on a useful
  timescale.
\item Deployed in chronic pain, neuropsychiatric PTSD, and rare-disease
  settings where parallel-group trials are infeasible.
\end{itemize}

\end{bullets}

\begin{rgt}
Lorem ipsum dolor sit amet, consectetur adipiscing elit, sed do eiusmod
tempor incididunt ut labore et dolore magna aliqua.

\end{rgt}

\begin{orig}
The aggregated N-of-1 trial design is a class of clinical trials in
which each enrolled participant cycles between active treatment and
control across multiple within-person blocks, with within-person
comparisons pooled across participants for population-level
inference\textsuperscript{\citeproc{ref-guyatt1986}{3}--\citeproc{ref-senn2018}{6}}.
The design preserves the within-subject contrast that gives
single-participant N-of-1 trials their inferential strength while
accumulating the statistical efficiency of a multi-participant trial.
For chronic conditions whose treatment response is reversible on a
useful timescale, the design yields more information per participant per
unit time than a parallel-group randomized trial of equivalent total
sample size\textsuperscript{\citeproc{ref-mirza2017}{7}}. It has been
deployed in chronic pain\textsuperscript{\citeproc{ref-yelland2009}{8}},
inflammatory bowel disease, neuropsychiatric symptoms in post-traumatic
stress disorder\textsuperscript{\citeproc{ref-hendrickson2020}{9}}, and
a growing portfolio of rare and orphan disease applications where
parallel-group trials are
infeasible\textsuperscript{\citeproc{ref-lillie2011}{5},\citeproc{ref-schork2015}{10}}.

\end{orig}

\begin{bullets}

\textbf{The primary scientific estimand in biomarker-guided N-of-1
trials is the biomarker-treatment interaction term.}

\begin{itemize}\setlength{\itemsep}{2pt}\setlength{\parskip}{0pt}
\item A predictive biomarker enables point-of-care treatment
  stratification per the PATH framework.
\item Power depends on moderation magnitude, within-participant
  contrast precision, trial design, and carryover leakage.
\item Carryover erodes the cleanness of within-subject contrasts and
  is the principal nuisance parameter in power calculations.
\end{itemize}

\end{bullets}

\begin{rgt}
Lorem ipsum dolor sit amet, consectetur adipiscing elit, sed do eiusmod
tempor incididunt ut labore et dolore magna aliqua.

\end{rgt}

\begin{orig}
Within this design class, a question of increasing scientific weight is
whether treatment response is moderated by a measured pre-treatment
biomarker. A biomarker that predicts which participants respond to which
treatments would allow clinicians to stratify treatment decisions at the
point of care, and the formal framework for such inference is laid out
in the PATH statement\textsuperscript{\citeproc{ref-kent2020path}{11}}.
The corresponding statistical estimand is the biomarker-treatment
interaction term in the within-participant outcome model. Power to
detect this interaction depends on the magnitude of moderation, the
precision of the within-participant contrasts, the trial design
(open-label, crossover, hybrid), and any time-dependent leakage between
treatment phases, conventionally termed `carryover', that erodes the
cleanness of the within-subject contrasts.

\end{orig}

\begin{bullets}

\textbf{The DGP encoding of biomarker moderation determines
carryover-induced power loss, a question the prior literature raised
but did not resolve.}

\begin{itemize}\setlength{\itemsep}{2pt}\setlength{\parskip}{0pt}
\item Power loss under carryover varies sharply with how moderation
  is encoded in the DGP.
\item Two formal answers from outside biostatistics -- continuous
  heterogeneity and latent-class structure -- dominate the available
  options.
\item This paper maps those traditions and their consequences for
  N-of-1 trial design and analysis.
\end{itemize}

\end{bullets}

\begin{rgt}
Lorem ipsum dolor sit amet, consectetur adipiscing elit, sed do eiusmod
tempor incididunt ut labore et dolore magna aliqua.

\end{rgt}

\begin{orig}
Recent methodological work has documented substantial power losses for
biomarker-treatment interaction inference under realistic carryover,
with the loss varying sharply by how the moderation effect is encoded in
the data-generating
process\textsuperscript{\citeproc{ref-hendrickson2020}{9}}. The present
paper takes up a question that the prior work raised but did not
resolve: how should biomarker moderation be represented in the DGP
itself? Two formal answers, drawn from a long history of work outside
biostatistics, dominate the available options. We begin by laying out
these two answers and the methodological tradition behind each, before
turning to the specific consequences for N-of-1 trial design and
analysis.

\end{orig}

\subsection{Two formal answers: continuous heterogeneity and
latent-class
structure}\label{two-formal-answers-continuous-heterogeneity-and-latent-class-structure}

\begin{bullets}

\textbf{The continuous-heterogeneity answer encodes biomarker
moderation as a smooth interaction in a single multivariate
distribution.}

\begin{itemize}\setlength{\itemsep}{2pt}\setlength{\parskip}{0pt}
\item Every participant has some treatment response; the response is a
  linear function of the biomarker.
\item Population heterogeneity is captured by a single distribution
  with an interaction parameter $\beta_{bm:D}$.
\item This is the dominant tradition in clinical biostatistics via
  random slopes and fixed-effects interaction terms.
\end{itemize}

\end{bullets}

\begin{rgt}
Lorem ipsum dolor sit amet, consectetur adipiscing elit, sed do eiusmod
tempor incididunt ut labore et dolore magna aliqua.

\end{rgt}

\begin{orig}
The first answer represents biomarker moderation as a graded continuous
effect: every participant has some degree of treatment response, the
response is a smooth (typically linear) function of the biomarker, and
population heterogeneity is captured by a single multivariate
distribution with an interaction parameter. This is the dominant
tradition in clinical biostatistics, where heterogeneity of treatment
effect is routinely modeled via random slopes on the treatment indicator
and via interaction terms in fixed-effects regression. In its most
economical form, the moderation enters the conditional mean of the
outcome as \(\beta_{bm:D} \cdot B \cdot D\), where \(B\) is the
biomarker and \(D\) is the (possibly continuous-decay) treatment
indicator.

\end{orig}

\begin{bullets}

\textbf{The latent-class answer encodes biomarker moderation as
discrete subpopulation structure, with the biomarker predicting class
membership.}

\begin{itemize}\setlength{\itemsep}{2pt}\setlength{\parskip}{0pt}
\item Participants belong to unobserved responder/non-responder classes
  with qualitatively distinct treatment-response curves.
\item Class membership is a Bernoulli draw with biomarker-dependent
  mixing probabilities.
\item Marginal moderation emerges from the mixture, not a smooth
  conditional-mean function.
\end{itemize}

\end{bullets}

\begin{rgt}
Lorem ipsum dolor sit amet, consectetur adipiscing elit, sed do eiusmod
tempor incididunt ut labore et dolore magna aliqua.

\end{rgt}

\begin{orig}
The second answer represents biomarker moderation as a discrete
subpopulation structure: participants belong to one of a small number of
unobserved classes (most commonly `responder' and `non-responder'), the
classes have qualitatively distinct treatment-response curves, and the
biomarker predicts class membership without determining it. The
corresponding generative form is a finite mixture in which class
membership is a Bernoulli or multinomial draw with mixing probabilities
that depend on the biomarker, and the within-class outcome model is
class-specific. The marginal moderation observed in the data emerges
from the mixture rather than from a smooth conditional-mean function.

\end{orig}

\begin{bullets}

\textbf{The two encodings reflect genuinely distinct biological
hypotheses that diverge under the longitudinal structure of an N-of-1
trial.}

\begin{itemize}\setlength{\itemsep}{2pt}\setlength{\parskip}{0pt}
\item Continuous heterogeneity implies a graded physiological substrate
  scaling drug effect; latent-class implies qualitatively distinct
  phenotypes.
\item Both views yield identical marginal associations at a single
  timepoint but diverge once carryover and serial dependence are
  engaged.
\item Methodology from outside biostatistics -- with mature
  formalization and identifiability analysis -- is directly applicable.
\end{itemize}

\end{bullets}

\begin{rgt}
Lorem ipsum dolor sit amet, consectetur adipiscing elit, sed do eiusmod
tempor incididunt ut labore et dolore magna aliqua.

\end{rgt}

\begin{orig}
These two answers correspond to genuinely different biological
hypotheses. The continuous-heterogeneity view is consistent with a
mechanism in which the biomarker reflects a graded physiological
substrate (receptor density, enzymatic activity, baseline trait
severity) that scales drug effect across the entire population. The
latent-class view is consistent with a mechanism in which qualitatively
distinct phenotypes exist (a metabolising versus non-metabolising
allelic variant; a present versus absent target receptor) and the
biomarker is an imperfect indicator of phenotype. The two views can
produce identical marginal biomarker-outcome associations at a single
timepoint and diverge sharply once the longitudinal structure of an
N-of-1 trial is engaged, because carryover and within-person serial
dependence interact differently with continuous and class-structured
heterogeneity. This is the substantive point at which methodology
imported from outside biostatistics becomes useful: both traditions have
long histories of formalization and identifiability analysis in their
respective fields, and the corresponding tools are mature.

\end{orig}

\subsection{Psychometric origins}\label{psychometric-origins}

\begin{bullets}

\textbf{The continuous-trait tradition traces to Spearman and the
factor-analytic lineage of Thurstone and Lawley, in which heterogeneity
is by assumption continuous.}

\begin{itemize}\setlength{\itemsep}{2pt}\setlength{\parskip}{0pt}
\item Observed indicators are modeled as linear combinations of
  unobserved continuous factors plus residual noise.
\item Between-person heterogeneity is summarized by factor scores along
  one or more dimensions.
\item This tradition forms the psychometric foundation for the
  continuous-heterogeneity encoding of biomarker moderation.
\end{itemize}

\end{bullets}

\begin{rgt}
Lorem ipsum dolor sit amet, consectetur adipiscing elit, sed do eiusmod
tempor incididunt ut labore et dolore magna aliqua.

\end{rgt}

\begin{orig}
The continuous-trait view in formal statistical work originates with
Spearman's two-factor theory of
intelligence\textsuperscript{\citeproc{ref-spearman1904}{12}}, in which
performance on cognitive tests was modeled as the sum of a single
general factor and test-specific factors. Spearman's machinery was
elaborated through the 1930s and 1940s by Thurstone, Hotelling, Lawley,
and others into the modern factor-analytic tradition, in which observed
indicator variables are written as linear combinations of unobserved
continuous factors plus residual
noise\textsuperscript{\citeproc{ref-thurstone1947}{13},\citeproc{ref-lawleymaxwell1971}{14}}.
Within this tradition, between-person heterogeneity in any latent
construct is by assumption continuous, and individual differences are
summarized by factor scores along one or more dimensions.

\end{orig}

\begin{bullets}

\textbf{The latent-class tradition originates with Lazarsfeld's local-
independence formulation and was placed on rigorous ML footing by
Goodman, then extended to longitudinal and covariate-dependent forms.}

\begin{itemize}\setlength{\itemsep}{2pt}\setlength{\parskip}{0pt}
\item Local independence within classes is the founding assumption of
  latent class analysis (LCA); Goodman derived identifiability
  conditions and ML estimation.
\item Extensions through the 1970s--1980s cover polytomous and
  continuous indicators, multilevel data, and covariate-dependent
  class probabilities.
\item McLachlan and Peel's monograph is the standard entry point for
  finite mixture models.
\end{itemize}

\end{bullets}

\begin{rgt}
Lorem ipsum dolor sit amet, consectetur adipiscing elit, sed do eiusmod
tempor incididunt ut labore et dolore magna aliqua.

\end{rgt}

\begin{orig}
The discrete-class view emerged from a parallel but methodologically
distinct literature beginning with Lazarsfeld's work on latent structure
analysis in the
1950s\textsuperscript{\citeproc{ref-lazarsfeld1950}{15},\citeproc{ref-lazarsfeldhenry1968}{16}}.
Lazarsfeld observed that in attitudinal and survey data, the joint
distribution of categorical responses to multiple items frequently
displayed structure inconsistent with a single continuous latent trait,
and instead suggested that respondents belonged to one of several
qualitatively distinct latent classes within which item responses were
locally independent. This `local independence' formulation is the
founding assumption of latent class analysis.
Goodman\textsuperscript{\citeproc{ref-goodman1974}{17}} placed the LCA
model on rigorous statistical footing, deriving identifiability
conditions and maximum-likelihood estimation procedures for the
binary-indicator case in a Biometrika paper that remains a standard
reference. Through the 1970s and 1980s, the LCA framework expanded to
polytomous and continuous indicators, the latter widely known as latent
profile analysis\textsuperscript{\citeproc{ref-mccutcheon1987}{18}}, to
multilevel and longitudinal
extensions\textsuperscript{\citeproc{ref-vermunt2003}{19}}, and to
mixtures with covariate-dependent class
probabilities\textsuperscript{\citeproc{ref-daytonmacready1988}{20},\citeproc{ref-bandeenroche1997}{21}}.
McLachlan and Peel\textsuperscript{\citeproc{ref-mclachlanpeel2000}{22}}
unified much of this material into a comprehensive monograph on finite
mixture models that is the standard entry point for quantitative
researchers from outside the discipline.

\end{orig}

\begin{bullets}

\textbf{The psychometric literature has never resolved the
continuous-versus-discrete question, motivating hybrid factor mixture
models.}

\begin{itemize}\setlength{\itemsep}{2pt}\setlength{\parskip}{0pt}
\item Bartholomew argued the data rarely settle the continuous-discrete
  choice without strong prior assumptions.
\item Cognitive ability favors continuous factors; certain
  psychopathologies favor discrete subtypes.
\item Yung and Arminger developed factor mixture models combining
  discrete classes with within-class continuous factors.
\end{itemize}

\end{bullets}

\begin{rgt}
Lorem ipsum dolor sit amet, consectetur adipiscing elit, sed do eiusmod
tempor incididunt ut labore et dolore magna aliqua.

\end{rgt}

\begin{orig}
The defining tension in psychometric methodology has been whether a
given pattern of observed heterogeneity is best understood as continuous
or as discrete.
Bartholomew\textsuperscript{\citeproc{ref-bartholomew1987}{23}} framed
this as the choice between `graded' and `unfolding' models of individual
differences and argued that the data themselves rarely settle the
question without strong prior assumptions. The empirical literature has
accumulated examples in both directions: cognitive ability is generally
well-described by continuous factor models, while certain forms of
psychopathology display latent-class structure consistent with
qualitatively distinct subtypes. Subsequent methodological work,
beginning with\textsuperscript{\citeproc{ref-yung1997}{24}}
and\textsuperscript{\citeproc{ref-arminger1999}{25}}, developed hybrid
frameworks (factor mixture models) in which discrete classes coexist
with within-class continuous factors. We return to these frameworks
below.

\end{orig}

\subsection{Econometric and marketing-science
development}\label{econometric-and-marketing-science-development}

\begin{bullets}

\textbf{Heckman and Singer's econometric work on unobserved
heterogeneity anticipated finite-mixture methodology, warning that
structural estimates are sensitive to the assumed heterogeneity
distribution.}

\begin{itemize}\setlength{\itemsep}{2pt}\setlength{\parskip}{0pt}
\item The econometric concern was unobserved heterogeneity in
  micro-econometric models where identical-agent assumptions fail.
\item Heckman and Singer proposed non-parametric ML with the
  heterogeneity distribution estimated as a discrete mixture.
\item Lindsay unified the theory of mixture likelihoods; this work
  directly informed biostatistical finite-mixture practice.
\end{itemize}

\end{bullets}

\begin{rgt}
Lorem ipsum dolor sit amet, consectetur adipiscing elit, sed do eiusmod
tempor incididunt ut labore et dolore magna aliqua.

\end{rgt}

\begin{orig}
A largely independent strand of work on finite mixtures developed in
econometrics and quantitative marketing through the 1980s and 1990s. The
motivating problems differed from psychometrics: rather than measurement
of cognitive or psychological constructs, the econometric concern was
unobserved heterogeneity in micro-econometric models (consumer choice,
labor-market participation, firm productivity) where standard regression
assumptions of identical agents were known to fail. Heckman and
Singer\textsuperscript{\citeproc{ref-heckmansinger1984}{26}}, in an
Econometrica paper on duration models, demonstrated that estimates of
structural parameters in survival-style models could be highly sensitive
to the assumed functional form of the unobserved heterogeneity
distribution, and recommended a non-parametric maximum-likelihood
approach in which the heterogeneity distribution itself was estimated as
a discrete mixture with unknown support points and weights. Their
estimator anticipated much of what would later become standard
finite-mixture methodology in biostatistics, and
Lindsay\textsuperscript{\citeproc{ref-lindsay1995}{27}} placed the
broader theory of mixture likelihoods on a unified mathematical footing.

\end{orig}

\begin{bullets}

\textbf{The marketing-science literature formalized covariate-dependent
gating functions and clusterwise regression, providing the vocabulary
now used in mixture-of-experts models.}

\begin{itemize}\setlength{\itemsep}{2pt}\setlength{\parskip}{0pt}
\item DeSarbo and Cron introduced clusterwise linear regression with
  cluster-specific coefficients; Wedel and DeSarbo extended it to GLMs.
\item The gating function -- class-membership probability as a function
  of covariates -- became a first-class object.
\item Jacobs et al. formalized the same structure as 'mixtures of
  experts' in machine learning, supplying the standard vocabulary.
\end{itemize}

\end{bullets}

\begin{rgt}
Lorem ipsum dolor sit amet, consectetur adipiscing elit, sed do eiusmod
tempor incididunt ut labore et dolore magna aliqua.

\end{rgt}

\begin{orig}
In quantitative marketing, the parallel concern was customer
segmentation. A population of consumers might respond to a marketing
intervention differently depending on which unobserved segment they
belonged to, and the firm needed both to estimate segment-specific
response functions and to predict segment membership for new consumers.
DeSarbo and Cron\textsuperscript{\citeproc{ref-desarbocron1988}{28}}
formalized this as `clusterwise linear regression', in which
observations were partitioned into clusters with cluster-specific
regression coefficients. Wedel and
DeSarbo\textsuperscript{\citeproc{ref-wedeldesarbo1995}{29}} generalized
the approach to generalized linear models, and Jedidi and
colleagues\textsuperscript{\citeproc{ref-jedidi1997}{30}} extended it to
structural equation models with unobserved segments. Within this
literature, the gating function (the model for the probability of
segment membership as a function of covariates) became a first-class
object, and methodological work focused on whether the gating model
could share covariates with the within-segment regression model. Jacobs
and colleagues\textsuperscript{\citeproc{ref-jacobs1991}{31}}, working
contemporaneously in machine learning, formalized an essentially
identical structure under the name `mixtures of experts', which provided
much of the now-standard vocabulary.

\end{orig}

\begin{bullets}

\textbf{Two econometric themes translate directly to the N-of-1 trial
setting: identifiability concerns at small $N$, and biomarker-as-gating-
covariate structure.}

\begin{itemize}\setlength{\itemsep}{2pt}\setlength{\parskip}{0pt}
\item Econometric identifiability concerns are acute when heterogeneity
  is the principal inference target rather than a nuisance.
\item The biomarker-treatment moderation question is whether the gating
  function depends on the biomarker and within-class slopes differ.
\item These two themes directly motivate the present paper's framing.
\end{itemize}

\end{bullets}

\begin{rgt}
Lorem ipsum dolor sit amet, consectetur adipiscing elit, sed do eiusmod
tempor incididunt ut labore et dolore magna aliqua.

\end{rgt}

\begin{orig}
Two themes from the econometric strand are particularly relevant to the
N-of-1 trial setting. First, the econometric tradition treated
unobserved heterogeneity as a nuisance in many applications, and was
correspondingly attentive to whether the heterogeneity distribution was
identified, weakly identified, or over-fit. The identifiability concerns
raised by\textsuperscript{\citeproc{ref-heckmansinger1984}{26}} and
elaborated by\textsuperscript{\citeproc{ref-lindsay1995}{27}} translate
directly into N-of-1 settings where sample sizes are small and the
heterogeneity is the principal object of inference rather than a
nuisance. Second, the marketing-science focus on covariate-dependent
gating functions corresponds exactly to the
biomarker-as-class-membership-predictor structure that motivates the
present paper. The biomarker-treatment moderation question, translated
into econometric vocabulary, is whether the gating function depends on a
single observed covariate (the biomarker) and whether the within- class
regression coefficient on treatment differs across classes.

\end{orig}

\subsection{Convergence on hybrid
frameworks}\label{convergence-on-hybrid-frameworks}

\begin{bullets}

\textbf{Psychometric and econometric traditions converged on factor
mixture models combining discrete class structure with within-class
continuous heterogeneity.}

\begin{itemize}\setlength{\itemsep}{2pt}\setlength{\parskip}{0pt}
\item Yung developed the factor mixture model (FMM) as a finite mixture
  of confirmatory factor
  models with class-specific loadings and residual covariances.
\item Arminger placed FMMs in the Mplus structural equation framework.
\item Lubke and Muthen provided the first systematic Monte Carlo
  evaluation of conditions under which class structure is
  distinguishable from continuous heterogeneity.
\end{itemize}

\end{bullets}

\begin{rgt}
Lorem ipsum dolor sit amet, consectetur adipiscing elit, sed do eiusmod
tempor incididunt ut labore et dolore magna aliqua.

\end{rgt}

\begin{orig}
By the late 1990s, both the psychometric and the econometric strands had
converged on a hybrid formulation: combine continuous within-class
factor structure with discrete between-class structure.
Yung\textsuperscript{\citeproc{ref-yung1997}{24}}, in Psychometrika,
developed the factor mixture model as a finite mixture of confirmatory
factor models, where each class had its own factor-loading and residual
covariance structure. Arminger and
colleagues\textsuperscript{\citeproc{ref-arminger1999}{25}} placed
factor mixture models in the broader Mplus-style structural equation
framework. Lubke and
Muthen\textsuperscript{\citeproc{ref-lubkemuthen2005}{2}} conducted the
first systematic Monte Carlo evaluation of factor mixture models and
clarified the conditions under which the discrete-class component was
empirically distinguishable from continuous heterogeneity. Clark and
colleagues\textsuperscript{\citeproc{ref-clark2013}{32}} surveyed the
modeling strategies and presented an applied analysis of psychiatric
comorbidity.

\end{orig}

\begin{bullets}

\textbf{Growth mixture models and group-based trajectory models extend
the finite-mixture framework to longitudinal outcome trajectories.}

\begin{itemize}\setlength{\itemsep}{2pt}\setlength{\parskip}{0pt}
\item Growth mixture models (GMMs) posit class-specific trajectories
  with within-class
  random-effects variation; group-based models constrain
  within-class variance to zero.
\item Nagin's group-based approach originated in criminology and is
  now widely used in epidemiology and clinical research.
\item Both families are the longitudinal analogues of the latent-class
  DGP developed in the present paper.
\end{itemize}

\end{bullets}

\begin{rgt}
Lorem ipsum dolor sit amet, consectetur adipiscing elit, sed do eiusmod
tempor incididunt ut labore et dolore magna aliqua.

\end{rgt}

\begin{orig}
In parallel, the longitudinal extension of these ideas yielded the
growth mixture
model\textsuperscript{\citeproc{ref-muthenshedden1999}{33}--\citeproc{ref-muthen2004}{35}}
and the closely related group-based trajectory
model\textsuperscript{\citeproc{ref-nagin1999}{36},\citeproc{ref-nagin2005}{37}}.
Both posit that a longitudinal outcome trajectory is generated by a
finite mixture of class-specific trajectories. Growth mixture models
permit within-class random-effects variation around the class-specific
mean trajectory; group-based trajectory models constrain within-class
variance to zero, so each class has a single deterministic trajectory.
Nagin's group-based approach was developed primarily in criminology to
characterize developmental trajectories of antisocial behavior, and has
been widely applied in epidemiology, public health, and clinical
research where trajectory heterogeneity is the principal scientific
question.

\end{orig}

\begin{bullets}

\textbf{Overextraction hazards tempered early optimizm: GMMs can
spuriously identify classes from non-gaussian single-component data,
and reliable class recovery requires large samples.}

\begin{itemize}\setlength{\itemsep}{2pt}\setlength{\parskip}{0pt}
\item Bauer and Curran showed GMMs extract apparent classes from
  non-gaussian single-component DGPs.
\item Lubke and Muthen documented analogous FMM hazards and identified
  joint conditions on $N$, class separation, and parameter differences.
\item Nylund et al. systematised class enumeration via Monte Carlo
  evaluation of information criteria and the bootstrap likelihood-ratio
  test (BLRT); the cumulative
  message is that hundreds to thousands of participants are required.
\end{itemize}

\end{bullets}

\begin{rgt}
Lorem ipsum dolor sit amet, consectetur adipiscing elit, sed do eiusmod
tempor incididunt ut labore et dolore magna aliqua.

\end{rgt}

\begin{orig}
The methodological optimizm around growth and factor mixture models
during the early 2000s was tempered by a series of papers documenting
overextraction hazards. Bauer and
Curran\textsuperscript{\citeproc{ref-bauercurran2003}{1}}, in
Psychological Methods, demonstrated through simulation that growth
mixture models fitted to data generated from a single non-gaussian
distribution can spuriously identify multiple latent classes; the
apparent classes are artifacts of distributional non-normality rather
than evidence of substantive subpopulation structure.
Bauer\textsuperscript{\citeproc{ref-bauer2007}{38}} elaborated on the
implications for substantive psychological research and recommended
caution in interpreting class structure without independent validation.
Lubke and Muthen\textsuperscript{\citeproc{ref-lubkemuthen2005}{2}}
documented the analogous concerns for factor mixture models, identifying
the joint conditions on sample size, class separation, and
class-specific parameter differences under which class structure is
recoverable. Nylund and
colleagues\textsuperscript{\citeproc{ref-nylund2007}{39}} systematised
the class-enumeration question through Monte Carlo evaluation of
information criteria and likelihood-ratio bootstrap procedures. The
cumulative message of this literature is that mixture models can extract
apparent class structure where none exists, and that distinguishing
genuine from spurious class structure requires either substantial
samples (typically several hundred to several thousand participants) or
strong prior information.

\end{orig}

\subsection{Translation into longitudinal
biostatistics}\label{translation-into-longitudinal-biostatistics}

\begin{bullets}

\textbf{Verbeke and Lesaffre introduced finite-mixture random-effects
distributions into biostatistics via longitudinal HIV data, with Zhang
and Davidian extending to semi-parametric specifications.}

\begin{itemize}\setlength{\itemsep}{2pt}\setlength{\parskip}{0pt}
\item The model augments a standard random-intercept-and-slope linear
  mixed model (LMM)
  with a discrete latent class indicator on the random-effects
  population.
\item Their motivating application was qualitatively distinct CD4
  trajectories in HIV-infected patient latent classes.
\item Zhang and Davidian allowed more flexible heterogeneity
  distributions via semi-parametric mixture specifications.
\end{itemize}

\end{bullets}

\begin{rgt}
Lorem ipsum dolor sit amet, consectetur adipiscing elit, sed do eiusmod
tempor incididunt ut labore et dolore magna aliqua.

\end{rgt}

\begin{orig}
Mixture and latent-class methods entered biostatistical practice through
several complementary channels. Verbeke and
Lesaffre\textsuperscript{\citeproc{ref-verbekelesaffre1996}{40}}, in the
Journal of the American Statistical Association, proposed a linear
mixed-effects model in which the random-effects population was itself a
finite mixture of multivariate normals. Their motivating application was
longitudinal HIV-infection data in which latent classes of patients
displayed qualitatively different CD4 trajectories. The model is
essentially a growth mixture model in biostatistical clothing, with the
standard random-intercept-and-slope linear mixed model augmented by a
discrete latent class indicator. Zhang and
Davidian\textsuperscript{\citeproc{ref-zhangdavidian2001}{41}} extended
the framework to allow more flexible random-effects distributions via
semi-parametric mixture specifications.

\end{orig}

\begin{bullets}

\textbf{The \texttt{lcmm} and \texttt{flexmix} R packages translate
finite-mixture longitudinal models into accessible software for
biostatistical applications.}

\begin{itemize}\setlength{\itemsep}{2pt}\setlength{\parskip}{0pt}
\item \texttt{lcmm} implements joint latent-class mixed models for
  longitudinal outcomes with optional joint survival modeling.
\item \texttt{flexmix} provides a general regression-mixture framework
  with user-definable component models including mixed-effects
  specifications.
\item Both have been used in pharmacokinetic, cognitive aging, and
  clinical latent-class discovery applications.
\end{itemize}

\end{bullets}

\begin{rgt}
Lorem ipsum dolor sit amet, consectetur adipiscing elit, sed do eiusmod
tempor incididunt ut labore et dolore magna aliqua.

\end{rgt}

\begin{orig}
Proust-Lima and colleagues developed a more comprehensive software
framework with the \texttt{lcmm} R
package\textsuperscript{\citeproc{ref-proustlima2017}{42}}, which
implements joint latent-class mixed models for longitudinal continuous
and ordinal outcomes, with optional joint modeling of survival outcomes
via shared latent class structure. The \texttt{lcmm} framework has been
deployed in cognitive aging research, Alzheimer's disease progression
modeling, and clinical-trial latent-class discovery. The
\texttt{flexmix} package in
R\textsuperscript{\citeproc{ref-leisch2004}{43},\citeproc{ref-grunleisch2008}{44}}
provides a more general infrastructure for fitting mixtures of
regression-style models with user-supplied component specifications, and
has been used in pharmacokinetic and pharmacodynamic modeling where
individual-level variation in dose-response curves motivates a mixture
structure.

\end{orig}

\begin{bullets}

\textbf{Mixture and latent-class methods have not been systematically
applied to aggregated N-of-1 designs; the biomarker-moderation
question has been encoded exclusively as a continuous interaction.}

\begin{itemize}\setlength{\itemsep}{2pt}\setlength{\parskip}{0pt}
\item Published N-of-1 methods focus on within-participant estimation,
  Bayesian hierarchical models, and meta-analytic aggregation.
\item Biomarker-moderation power calculations use continuous
  interaction terms with no parallel latent-class formulation.
\item The mixture-formulation question appears unaddressed in the
  existing N-of-1 methodology literature.
\end{itemize}

\end{bullets}

\begin{rgt}
Lorem ipsum dolor sit amet, consectetur adipiscing elit, sed do eiusmod
tempor incididunt ut labore et dolore magna aliqua.

\end{rgt}

\begin{orig}
Despite this translation activity, mixture and latent-class methods have
not, to the best of our knowledge, been systematically applied to
aggregated N-of-1 trial designs. The published N-of-1 methods literature
is dominated by within-participant effect
estimation\textsuperscript{\citeproc{ref-lillie2011}{5},\citeproc{ref-zucker1997}{45}},
Bayesian hierarchical models for participant-specific
effects\textsuperscript{\citeproc{ref-zucker2010}{46}}, and
aggregated-evidence meta-analytic
frameworks\textsuperscript{\citeproc{ref-senn2018}{6},\citeproc{ref-schork2015}{10}}.
The biomarker-moderation question has received occasional treatment
under the heading of precision-medicine power
calculation\textsuperscript{\citeproc{ref-hendrickson2020}{9},\citeproc{ref-gabler2011}{47}},
but in each case the moderation has been encoded as a continuous
interaction term in a fixed-effects analysis, with no parallel treatment
of the latent-class formulation. The mixture-formulation question
appears to be unaddressed in the existing N-of-1 methodology literature.

\end{orig}

\subsection{Why this matters for N-of-1 designs
specifically}\label{why-this-matters-for-n-of-1-designs-specifically}

Three features of the aggregated N-of-1 design make the
continuous-versus-discrete moderation question particularly
consequential for design and analysis decisions.

\begin{bullets}

\textbf{Under a latent-class DGP, class membership is fixed across
drug states and is in principle recoverable from the full trajectory
even when on-drug contrasts are noisy.}

\begin{itemize}\setlength{\itemsep}{2pt}\setlength{\parskip}{0pt}
\item Continuous-heterogeneity DGPs encode moderation in how the
  within-participant contrast varies with the biomarker.
\item Latent-class DGPs encode moderation in class membership, which
  is invariant to drug state.
\item The two encodings imply different efficient analysis strategies
  and different power profiles under carryover.
\end{itemize}

\end{bullets}

\begin{rgt}
Lorem ipsum dolor sit amet, consectetur adipiscing elit, sed do eiusmod
tempor incididunt ut labore et dolore magna aliqua.

\end{rgt}

\begin{orig}
First, the within-participant longitudinal structure of an N-of-1 trial
provides repeated measurements at multiple drug states (on-drug,
off-drug, in carryover). Under a continuous-heterogeneity DGP, the
moderation signal is encoded in how the within-participant treatment
contrast varies with the biomarker. Under a latent-class DGP, the
moderation signal is encoded in class membership, which is fixed across
all measurements within a participant and is therefore in principle
recoverable from the full trajectory even when individual on-drug
contrasts are noisy. The two encodings imply different efficient
analysis strategies and different power profiles under
carryover\textsuperscript{\citeproc{ref-hendrickson2020}{9}}.

\end{orig}

\begin{bullets}

\textbf{N-of-1 trial sample sizes are well below the threshold at which
mixture-model identifiability has been validated, so a misspecified
single-component analysis may be the more reliable tool.}

\begin{itemize}\setlength{\itemsep}{2pt}\setlength{\parskip}{0pt}
\item Psychometric identifiability validation requires hundreds to
  thousands of participants; N-of-1 trials operate at $N \in [30,150]$.
\item Even a biologically correct latent-class DGP may yield
  underpowered or unidentified mixture estimates at trial-relevant $N$.
\item The mixture-versus-mitigation trade-off is an empirical question
  not addressed at N-of-1-relevant sample sizes.
\end{itemize}

\end{bullets}

\begin{rgt}
Lorem ipsum dolor sit amet, consectetur adipiscing elit, sed do eiusmod
tempor incididunt ut labore et dolore magna aliqua.

\end{rgt}

\begin{orig}
Second, N-of-1 trials are typically conducted at sample sizes of
\(N \in [30, 150]\), well below the sample sizes at which mixture-model
identifiability has been validated in the psychometric literature. The
identifiability cautions documented
by\textsuperscript{\citeproc{ref-bauercurran2003}{1}}
and\textsuperscript{\citeproc{ref-lubkemuthen2005}{2}} therefore apply
with full force, and the analysis-strategy implications are not trivial:
even if a latent-class DGP is the biologically correct generative model,
mixture analysis at trial-relevant sample sizes may be either
underpowered or unidentified, and a misspecified single-component
analysis may be the more reliable inferential tool despite its bias.
Resolving this trade-off is in part an empirical question that the
existing psychometric simulation literature has not addressed at
N-of-1-relevant sample sizes.

\end{orig}

\begin{bullets}

\textbf{Carryover erodes the moderation signal directly under
continuous-heterogeneity but only indirectly under latent-class,
predicting divergent power-loss profiles.}

\begin{itemize}\setlength{\itemsep}{2pt}\setlength{\parskip}{0pt}
\item Under continuous heterogeneity, carryover attenuates $D_{bc,it}$
  and thereby directly attenuates the moderation signal.
\item Under latent-class structure, carryover erodes within-class drug
  contrasts but does not directly attenuate class membership.
\item The two formulations predict different power-loss profiles; this
  difference is empirically large in the companion carryover-
  sensitivity work.
\end{itemize}

\end{bullets}

\begin{rgt}
Lorem ipsum dolor sit amet, consectetur adipiscing elit, sed do eiusmod
tempor incididunt ut labore et dolore magna aliqua.

\end{rgt}

\begin{orig}
Third, carryover, which enters the analysis as a continuous decay in the
effective drug indicator \(D_{bc,it} = \exp(-\lambda t_{sd,it})\),
interacts differently with the two moderation encodings. In the
continuous-heterogeneity formulation, carryover erodes the moderation
signal directly because the moderation enters as a product with
\(D_{bc,it}\). In the latent-class formulation, carryover erodes the
within-class drug contrast but does not directly attenuate class
membership; if class membership can be recovered from the full
trajectory, some moderation signal can survive carryover that would be
lost under the continuous-heterogeneity analysis. The two formulations
therefore predict different power-loss profiles under carryover, and
this difference is empirically large in the simulation literature
accumulated by the companion carryover-sensitivity work.

\end{orig}

\subsection{Aims and structure of the
paper}\label{aims-and-structure-of-the-paper}

\begin{bullets}

\textbf{The paper has four aims: formalize latent-class DGPs for N-of-1
trials, characterize estimation strategies at small $N$, benchmark
power via simulation, and re-analyze the prazosin-PTSD dataset.}

\begin{itemize}\setlength{\itemsep}{2pt}\setlength{\parskip}{0pt}
\item Aims 1--2 are methodological: formal specification and
  estimation/diagnostic characterization with explicit identifiability
  conditions.
\item Aim 3 is computational: a factorial simulation benchmarking
  mixture-DGP against continuous-DGP power under carryover.
\item Aim 4 is applied: re-analysis of the Hendrickson et al. (2020)
  prazosin-PTSD data under the latent-class formulation.
\end{itemize}

\end{bullets}

\begin{rgt}
Lorem ipsum dolor sit amet, consectetur adipiscing elit, sed do eiusmod
tempor incididunt ut labore et dolore magna aliqua.

\end{rgt}

\begin{orig}
This paper has four aims. The first is to formalize the latent-class and
finite-mixture formulations of biomarker-treatment interaction in
aggregated N-of-1 trials, with explicit identifiability conditions
specific to the trial-relevant sample-size regime. The second is to
characterize estimation strategies (expectation-maximization, Bayesian
samplers, candidate R packages) and the diagnostics required to
distinguish genuine from spurious class structure at small \(N\). The
third is to report a factorial simulation study that benchmarks
mixture-DGP and continuous-DGP power against the standard
linear-mixed-model analysis and a class-aware analysis under realistic
carryover. The fourth is to apply the developed methods to the
prazosin-PTSD biomarker dataset
of\textsuperscript{\citeproc{ref-hendrickson2020}{9}} and to compare
conclusions with those obtained under the continuous-moderation analysis
reported in the original paper.

\end{orig}

\begin{bullets}

\textbf{The paper moves from notation and DGP specification through
multivariate normal (MVN) approximation, identifiability analysis,
model taxonomy, and spectrum
organization to analysis implications and a discussion of open
questions.}

\begin{itemize}\setlength{\itemsep}{2pt}\setlength{\parskip}{0pt}
\item Sections on the faithful DGP, MVN approximation, and failure
  regimes develop the theoretical core.
\item The finite-mixture taxonomy and biomarker-moderation spectrum
  organize the model landscape systematically.
\item The discussion addresses three open questions and situates the
  companion five-study simulation program.
\end{itemize}

\end{bullets}

\begin{rgt}
Lorem ipsum dolor sit amet, consectetur adipiscing elit, sed do eiusmod
tempor incididunt ut labore et dolore magna aliqua.

\end{rgt}

\begin{orig}
The remainder of the paper is organized as follows. We begin in the
notation section with a glossary that fixes terminology used throughout.
The latent-class generative form, its MVN second-moment approximation,
and the regimes in which the approximation breaks down are developed in
three subsequent sections. The implications for power under carryover
and the identifiability constraints that bound mixture analysis at
trial-relevant sample sizes are addressed next. A finite-mixture
taxonomy section locates the relevant model families (LCA, GMM, FMM,
regression mixtures); a biomarker-moderation spectrum section organizes
the available DGP encodings moment by moment. Implications for analysis
and design, and software for mixture analyses, are then drawn together.
The discussion section addresses three open methodological questions and
references the five-study simulation program described in the companion
simulation plan, which constitutes the empirical core of the present
line of work and is the next deliverable in the program.

\end{orig}

\section{Notation}\label{notation}

\begin{bullets}

\textbf{Core notation: $B_i$ is the biomarker, $BR_{it}$ the
biological-response component, $D_{bc,it}$ the continuous-decay drug
indicator, and $Z_i$ the unobserved class label.}

\begin{itemize}\setlength{\itemsep}{2pt}\setlength{\parskip}{0pt}
\item Outcome is decomposed into $BR_{it}$, $TV_{it}$, and $PB_{it}$
  following the Hendrickson framework.
\item $D_{bc,it} = \exp(-\lambda t_{sd,it})$ off-drug and $1$ on-drug;
  $\lambda = \ln 2 / t_{1/2}$.
\item The DGP and analysis model are permitted to use different
  carryover half-lives.
\end{itemize}

\end{bullets}

\begin{rgt}
Lorem ipsum dolor sit amet, consectetur adipiscing elit, sed do eiusmod
tempor incididunt ut labore et dolore magna aliqua.

\end{rgt}

\begin{orig}
Throughout, \(B_i\) denotes the pre-treatment biomarker value for
participant \(i\). \(BR_{it}\) denotes the biological-response component
of outcome for participant \(i\) at time \(t\), decomposed in the
Hendrickson framework into three latent response factors: \(BR_{it}\)
itself, a time-variant component \(TV_{it}\), and a placebo-belief
component \(PB_{it}\). \(D_{bc,it}\) denotes drug state at time \(t\) in
the analysis-side continuous-decay form
\(D_{bc,it} = \exp(-\lambda t_{sd,it})\) during off-drug periods and
\(D_{bc,it} = 1\) during on-drug periods, with \(t_{sd,it}\) the time
since drug discontinuation and \(\lambda\) the decay rate. Following the
convention used throughout this series, \(D_{it}\) is reserved for the
binary on-drug state and \(D_{bc,it}\) for the continuous
exposure-decayed indicator used here; the latter is the \texttt{Dbc}
regressor of the analysis model. \(Z_i\) denotes an unobserved class
label. The half-life of carryover is \(t_{1/2} = \ln 2 / \lambda\) and
is allowed to differ between the DGP and the analysis model.

\end{orig}

The principal psychometric and econometric terms used below are:

\begin{itemize}
\tightlist
\item
  \textbf{Latent class.} An unobserved categorical class, equivalent in
  trial vocabulary to an unobserved responder/non-responder class.
\item
  \textbf{Latent factor.} An unobserved continuous trait governing
  covariation among observed variables, equivalent to a frailty (the
  survival-analysis term for an unobserved multiplicative random effect)
  or to a random intercept in a linear mixed model.
\item
  \textbf{Class-conditional distribution \(f_z\).} The outcome
  distribution within a single latent class.
\item
  \textbf{Mixing proportion \(\pi\).} The prevalence of a given latent
  class in the population.
\item
  \textbf{Measurement invariance.} Whether a model applies identically
  across subpopulations.
\item
  \textbf{Growth mixture model (GMM).} A longitudinal mixed model in
  which the population of random effects is itself a finite mixture.
\item
  \textbf{Factor mixture model (FMM).} A hybrid model with discrete
  classes at the upper level and within-class continuous heterogeneity
  at the lower level.
\item
  \textbf{Regression mixture / mixture of experts.} A finite mixture in
  which the gating function (class-membership probability) depends on
  covariates and within-class regressions are class-specific.
\item
  \textbf{EM algorithm.} Expectation-maximization (EM) applied to
  recover unobserved class labels jointly with class-conditional model
  parameters.
\end{itemize}

\section{Faithful latent-class data-generating
processes}\label{faithful-latent-class-data-generating-processes}

If a candidate predictive biomarker indexes an unobserved
responder/non-responder dichotomy, the mechanistically faithful
data-generating process is a finite mixture:

\[
Z_i \mid B_i \sim \mathrm{Bernoulli}\bigl(\pi(B_i)\bigr), \qquad
BR_{it} \mid Z_i = z \sim f_z(t, D_{bc,it}),
\]

where \(Z_i\) is the latent class indicator, \(\pi(B_i) = \Pr(Z_i = 1
\mid B_i)\) is the class-membership probability (typically a logistic
function of the biomarker), and the class-conditional response
distributions \(f_0\) and \(f_1\) differ in their drug-response
parameters: peak response, onset rate, or both. Within-class temporal
structure follows the same modGompertz form used in single-class
formulations, with class-specific parameters substituted.

\begin{bullets}

\textbf{The faithful latent-class DGP is a finite mixture with
biomarker-dependent class-membership probability and class-specific
modGompertz drug-response distributions.}

\begin{itemize}\setlength{\itemsep}{2pt}\setlength{\parskip}{0pt}
\item $Z_i \mid B_i \sim \mathrm{Bernoulli}(\pi(B_i))$ with
  $\pi(B_i)$ typically a logistic function of the biomarker.
\item Class-conditional distributions $f_0$ and $f_1$ differ in
  drug-response parameters (peak response, onset rate, or both).
\item Within-class temporal structure uses the same modGompertz form
  as single-class formulations.
\end{itemize}

\end{bullets}

\begin{rgt}
Lorem ipsum dolor sit amet, consectetur adipiscing elit, sed do eiusmod
tempor incididunt ut labore et dolore magna aliqua.

\end{rgt}

\begin{orig}
If a candidate predictive biomarker indexes an unobserved
responder/non-responder dichotomy, the mechanistically faithful
data-generating process is a finite mixture, with \(Z_i \mid B_i\)
following a Bernoulli distribution with class-membership probability
\(\pi(B_i)\) (typically logistic in the biomarker), and the
class-conditional response distributions \(f_0\) and \(f_1\) differing
in their drug-response parameters. Within-class temporal structure uses
the same modGompertz form as single-class formulations with
class-specific parameters substituted.

\end{orig}

\begin{bullets}

\textbf{Class-label invariance preserves moderation information
across carryover that the MVN-covariance encoding, where the signal is
entirely in the second moment, cannot retain.}

\begin{itemize}\setlength{\itemsep}{2pt}\setlength{\parskip}{0pt}
\item $Z_i$ is fixed across drug states; it carries participant-level
  moderation information even when between-class mean separation is
  attenuated.
\item Under a full carryover off-drug period, class-conditional means
  converge
  but $Z_i$ remains recoverable from the full trajectory.
\item Under MVN-covariance encoding the carryover erosion of the
  covariance is direct and no residual signal survives the off-drug
  period.
\end{itemize}

\end{bullets}

\begin{rgt}
Lorem ipsum dolor sit amet, consectetur adipiscing elit, sed do eiusmod
tempor incididunt ut labore et dolore magna aliqua.

\end{rgt}

\begin{orig}
We note two consequences of this DGP form that are central to what
follows. First, the class label \(Z_i\) is fixed across all measurements
within a participant; it is invariant to drug state and to time. Second,
the class-conditional mean trajectory does depend on \(D_{bc,it}\):
under a full carryover off-drug period (\(D_{bc,it} \to 0\)) the two
class-conditional means converge if the responder/non-responder
distinction is purely in drug response. Class-label invariance therefore
preserves participant-level information about the moderation signal that
an analysis with access to \(Z_i\) could exploit, even where
instantaneous between-class mean separation in \(BR\) is fully
attenuated. The MVN-covariance encoding of biomarker moderation, in
contrast, has no such latent-label information to exploit: the
covariance is the entire signal, and its erosion under carryover is
direct.

\end{orig}

\section{MVN approximation to mixture-induced
moderation}\label{mvn-approximation-to-mixture-induced-moderation}

A two-component mixture with class-dependent means and common
within-class covariance produces a marginal joint distribution of
\((B_i, BR_{it})\) whose first two moments are reproduced, to leading
order, by a single multivariate normal with \(\mathrm{Cor}(B, BR) =
c_{bm}\). Under within-class independence between \(B_i\) and
\(BR_{it}\),

\[
c_{bm}^2 \;=\;
\frac{\pi(1-\pi)(\mu_B^1-\mu_B^0)^2}{\mathrm{Var}(B)}
\;\cdot\;
\frac{\pi(1-\pi)(\mu_{BR}^1-\mu_{BR}^0)^2}{\mathrm{Var}(BR)},
\]

so that \(c_{bm}^2\) is the product of the between-class variance
fractions in \(B\) and in \(BR\). The marginal correlation that an MVN
can reproduce is therefore bounded by the geometric mean of the two
intra-class correlations; if either class separation is small, the
achievable marginal \(c_{bm}\) is small even when class structure is
otherwise well-defined.

\begin{bullets}

\textbf{The achievable marginal $c_{bm}$ under an MVN approximation
is bounded by the geometric mean of the two between-class variance
fractions.}

\begin{itemize}\setlength{\itemsep}{2pt}\setlength{\parskip}{0pt}
\item $c_{bm}^2$ equals the product of the between-class variance
  fractions in $B$ and in $BR$.
\item Small class separation in either variable limits the achievable
  marginal $c_{bm}$ even when class structure is otherwise
  well-defined.
\item The MVN approximation reproduces the mixture's first two moments
  under within-class independence.
\end{itemize}

\end{bullets}

\begin{rgt}
Lorem ipsum dolor sit amet, consectetur adipiscing elit, sed do eiusmod
tempor incididunt ut labore et dolore magna aliqua.

\end{rgt}

\begin{orig}
A two-component mixture with class-dependent means and common
within-class covariance produces a marginal joint distribution of
\((B_i, BR_{it})\) whose first two moments are reproduced, to leading
order, by a single multivariate normal with \(\mathrm{Cor}(B, BR) =
c_{bm}\). The achievable marginal \(c_{bm}\) is bounded by the geometric
mean of the two between-class variance fractions; if either class
separation is small, the achievable \(c_{bm}\) is small even when class
structure is otherwise well-defined.

\end{orig}

\begin{bullets}

\textbf{Within the regularity conditions, the MVN second-moment
approximation is adequate for LMM power calculations and captures
differential correlation across drug states.}

\begin{itemize}\setlength{\itemsep}{2pt}\setlength{\parskip}{0pt}
\item The approximation yields elevated correlation on-drug and
  shrinkage toward zero off-drug after carryover has decayed.
\item Adequacy requires adequate class overlap and $\pi$ bounded
  away from 0 and 1.
\item For $t$-tests on $\beta_{bm:D}$ this approximation captures
  covariance consequences without committing to a discrete DGP.
\end{itemize}

\end{bullets}

\begin{rgt}
Lorem ipsum dolor sit amet, consectetur adipiscing elit, sed do eiusmod
tempor incididunt ut labore et dolore magna aliqua.

\end{rgt}

\begin{orig}
Under mild regularity conditions (adequate class overlap and mixing
proportions \(\pi\) bounded away from 0 and 1), the mixture's
second-moment structure is well-approximated by a single MVN with
differential correlation across drug states: elevated correlation where
class-conditional means separate (on-drug) and shrinkage towards zero
where they coincide (off-drug after carryover has decayed). For power
calculations driven by t-tests on \(\beta_{bm:D}\) in a linear
mixed-effects model, this approximation is adequate within the
regularity conditions and captures the covariance consequences of
latent-class structure without committing to a discrete generative form.

\end{orig}

\section{Failure regimes of the MVN
approximation}\label{failure-regimes-of-the-mvn-approximation}

The MVN approximation is misleading in three identifiable regimes, each
with a distinct empirical signature.

\begin{itemize}
\tightlist
\item
  \textbf{Bimodal responses.} When \(f_0\) and \(f_1\) are
  well-separated (non-responders show near-zero drug effect and
  responders show a large effect), the marginal distribution of \(BR\)
  is bimodal: a histogram of on-drug response would show two peaks
  rather than a single peak with gaussian scatter around it. No MVN
  matches the tails of a bimodal distribution, and power under the true
  mixture DGP will exceed the MVN-LMM prediction because a
  well-specified mixture analysis can exploit the bimodality. The
  pattern occurs in clinical settings where a subpopulation genuinely
  does not respond to a treatment, as in certain targeted oncology
  indications.
\item
  \textbf{Strong class-membership gradient.} When \(\pi(B_i)\) is steep,
  approaching a step function at a biomarker threshold, the biomarker
  behaves nearly deterministically as a class label. A thresholded
  mean-moderation specification (a thresholded Architecture-A variant)
  then becomes the better approximation, because the biomarker-response
  relationship is dominated by its mean-structure component. The
  corresponding clinical scenario is a predictive biomarker with
  near-perfect sensitivity and specificity for the responsive phenotype.
\item
  \textbf{Class-varying covariance.} If autocorrelation, residual
  variance, or placebo-response parameters differ between classes, no
  single-component MVN can represent the generative structure. Analysis
  models assuming constant covariance are misspecified regardless of
  whether they encode the biomarker interaction in the mean or in the
  covariance. This is the latent-class analogue of heteroscedastic
  residual variance addressed in standard mixed- effects practice with
  \texttt{varIdent} or \texttt{weights} specifications, except that the
  stratifying variable is unobserved.
\end{itemize}

\section{Implications for power under
carryover}\label{implications-for-power-under-carryover}

\begin{bullets}

\textbf{The power gap between a second-moment analysis and a
mixture-faithful analysis upper-bounds what a class-aware procedure
can recover, but the bound is asymptotic and may reverse at small $N$.}

\begin{itemize}\setlength{\itemsep}{2pt}\setlength{\parskip}{0pt}
\item The 40--60\% covariance-only power loss reflects direct erosion
  of a second-moment signal; class membership is carryover-invariant.
\item The gap between covariance-only and mixture-faithful power
  estimates upper-bounds the recoverable power.
\item At finite $N$ in the trial-relevant range, finite-sample bias
  can drive the realized gap below the bound or reverse its sign.
\end{itemize}

\end{bullets}

\begin{rgt}
Lorem ipsum dolor sit amet, consectetur adipiscing elit, sed do eiusmod
tempor incididunt ut labore et dolore magna aliqua.

\end{rgt}

\begin{orig}
The substantial covariance-only power loss under carryover (40--60\%
across designs in\textsuperscript{\citeproc{ref-hendrickson2020}{9}})
reflects erosion of a second-moment signal. A mixture-model analysis
fitted to mixture-DGP data can in principle recover additional power by
estimating class membership directly, because the class label is
carryover-invariant. The gap between the covariance-only
carryover-sensitive power estimate and the mixture-faithful power
estimate upper-bounds what a mixture analysis could recover; whether
this bound is tight in the trial-relevant parameter regime is an
empirical question not addressed by existing simulation infrastructure.
The bound is asymptotic: at finite \(N\), finite-sample bias can drive
the realized gap below the bound or, in adverse identifiability regimes,
reverse its sign.

\end{orig}

\begin{bullets}

\textbf{The Study 1 factorial crosses two DGPs and two analyses to
decompose approximation error from analysis-model mis-specification
and to verify Type I error control.}

\begin{itemize}\setlength{\itemsep}{2pt}\setlength{\parskip}{0pt}
\item The mixture-DGP/MVN-analysis vs.\ mixture-DGP/mixture-analysis
  contrast quantifies power recoverable by a faithful analysis.
\item The two MVN-analysis cells quantify approximation error in the
  second-moment encoding.
\item A null check on the MVN-DGP/mixture-analysis cell verifies no
  spurious power extraction.
\end{itemize}

\end{bullets}

\begin{rgt}
Lorem ipsum dolor sit amet, consectetur adipiscing elit, sed do eiusmod
tempor incididunt ut labore et dolore magna aliqua.

\end{rgt}

\begin{orig}
A focused factorial simulation isolating second-moment approximation
error from analysis-model mis-specification is described in the
companion simulation plan (Study 1: mixture-vs-MVN factorial). The
factorial crosses two DGPs (mixture and MVN) with two analyses (linear
mixed model and joint latent-class mixed model). The
mixture-DGP/MVN-analysis versus mixture-DGP/mixture-analysis contrast
quantifies power recoverable by a faithful analysis model; the
difference between the two MVN-analysis cells quantifies approximation
error. A null check on the MVN-DGP/mixture-analysis cell verifies the
mixture model does not extract spurious power on data with no latent
class structure.

\end{orig}

\section{Identifiability at trial-relevant sample
sizes}\label{identifiability-at-trial-relevant-sample-sizes}

\begin{bullets}

\textbf{Mixture EM requires either informative class-probability priors
or strong class separation; neither is assured at N-of-1 sample sizes,
which fall well below identifiability validation thresholds.}

\begin{itemize}\setlength{\itemsep}{2pt}\setlength{\parskip}{0pt}
\item Posterior entropy stabilization requires hundreds to thousands
  of participants in published psychometric applications.
\item N-of-1 trials operate at $N \in [30, 150]$, well below this
  range.
\item The identifiability cost is acute precisely when the moderation
  effect's existence is uncertain.
\end{itemize}

\end{bullets}

\begin{rgt}
Lorem ipsum dolor sit amet, consectetur adipiscing elit, sed do eiusmod
tempor incididunt ut labore et dolore magna aliqua.

\end{rgt}

\begin{orig}
Mixture modeling is attractive on faithfulness grounds but carries
identifiability costs that are acute in N-of-1 settings. The EM
algorithm for finite mixtures of mixed-effects models requires either
informative priors on class-membership probabilities or strong
separation between class-conditional response curves; neither is assured
in a trial designed to detect a moderation effect whose existence is
itself uncertain. The standard diagnostic for mixture identifiability,
the entropy of posterior class-membership probabilities, requires
substantial sample sizes (\(N\) several hundred to several thousand) to
stabilize at values that support confident class extraction in published
psychometric
applications\textsuperscript{\citeproc{ref-bauercurran2003}{1},\citeproc{ref-bauer2007}{38},\citeproc{ref-nylund2007}{39}}.
Aggregated N-of-1 trials of the size contemplated in current biomarker
applications, \(N \in [30, 150]\), are well below this range.

\end{orig}

\begin{bullets}

\textbf{The operative question is not faithfulness but whether the
fidelity gain is recoverable at realistic $N$, or whether standard
analysis mitigations offer a better power-per-complexity trade.}

\begin{itemize}\setlength{\itemsep}{2pt}\setlength{\parskip}{0pt}
\item Mixture analysis may be underpowered or unidentified even under
  a genuinely mixture-structured DGP at trial-relevant $N$.
\item Analysis-strategy mitigations within standard mixed-effects
  practice have better-understood inferential properties at small $N$.
\item Study 5 of the companion simulation plan addresses this
  directly with small-$N$ identifiability and Type I error checks.
\end{itemize}

\end{bullets}

\begin{rgt}
Lorem ipsum dolor sit amet, consectetur adipiscing elit, sed do eiusmod
tempor incididunt ut labore et dolore magna aliqua.

\end{rgt}

\begin{orig}
The methodological implication is that mixture analysis at
trial-relevant sample sizes may be either underpowered or unidentified,
even when the underlying DGP is genuinely mixture-structured. The
question is therefore not whether mixture formulations are more
faithful, but whether the fidelity gain is recoverable at realistic
sample sizes, or whether analysis-strategy mitigations available within
standard mixed-effects practice (restricted analysis, weighting,
within-subject contrasts;
see\textsuperscript{\citeproc{ref-hendrickson2020}{9}}) deliver a better
power-per-complexity trade. These mitigations have better-understood
inferential properties at the relevant sample sizes. Resolving the
mixture-versus-mitigation question empirically requires the simulation
program described in the companion simulation plan; Study 5 in
particular addresses small-\(N\) identifiability and Type I error
explicitly.

\end{orig}

\section{Finite-mixture taxonomy}\label{finite-mixture-taxonomy}

Four families of finite mixture model from the psychometric and
econometric literatures bear directly on the biomarker-treatment
interaction problem. In this section we shall introduce each family, its
origin, and its specific relevance to the N-of-1 trial setting.

\subsection{Latent class and latent profile
analysis}\label{latent-class-and-latent-profile-analysis}

\begin{bullets}

\textbf{LCA and latent profile analysis (LPA) are cross-sectional
mixture models estimating
class membership from observed indicators; in the trial context they
characterize responder classes from response trajectories.}

\begin{itemize}\setlength{\itemsep}{2pt}\setlength{\parskip}{0pt}
\item LCA uses categorical indicators; LPA uses continuous indicators;
  both return posterior class-membership probabilities.
\item In the biomarker-trial context, an LPA-style model estimates
  responder-class probabilities from on-drug and off-drug response
  trajectories.
\item Classical estimation is via EM; standard R implementations
  include \texttt{poLCA} and Latent GOLD.
\end{itemize}

\end{bullets}

\begin{rgt}
Lorem ipsum dolor sit amet, consectetur adipiscing elit, sed do eiusmod
tempor incididunt ut labore et dolore magna aliqua.

\end{rgt}

\begin{orig}
Latent class analysis and latent profile analysis are mixture models for
cross-sectional data. Each participant belongs to one of a small number
of unobserved classes; the observed variables have class-specific
marginal distributions. LCA uses categorical indicators; LPA uses
continuous indicators. A mixture fitted to baseline biomarkers alone
returns posterior class-membership probabilities that can be used as
covariates in subsequent outcome analysis, the two-stage analytic
strategy standard in psychometric applications. In the biomarker-trial
context, an LPA-style model posits two or more responder classes with
class-specific marginal distributions of on-drug and off-drug response,
and estimates class-membership probabilities from the observed response
trajectories rather than from baseline covariates alone. Classical
references
include\textsuperscript{\citeproc{ref-lazarsfeldhenry1968}{16}},\textsuperscript{\citeproc{ref-goodman1974}{17}},\textsuperscript{\citeproc{ref-clogg1995}{48}},
and\textsuperscript{\citeproc{ref-mclachlanpeel2000}{22}}; the Vermunt
and Magidson Latent GOLD software and the \texttt{poLCA} R package
implement the standard EM estimation routines.

\end{orig}

\subsection{Growth mixture models and group-based trajectory
models}\label{growth-mixture-models-and-group-based-trajectory-models}

\begin{bullets}

\textbf{GMMs extend LPA to longitudinal data with class-specific
drug-response curves and within-class random-effects variation;
group-based trajectory models constrain within-class variance to zero.}

\begin{itemize}\setlength{\itemsep}{2pt}\setlength{\parskip}{0pt}
\item A GMM is a longitudinal LMM with the random-effects population
  replaced by a mixture of multivariate normals.
\item Group-based trajectory models are the degenerate case with zero
  within-class random-effect variance.
\item Both families model longitudinal trajectory heterogeneity as
  the principal scientific question.
\end{itemize}

\end{bullets}

\begin{rgt}
Lorem ipsum dolor sit amet, consectetur adipiscing elit, sed do eiusmod
tempor incididunt ut labore et dolore magna aliqua.

\end{rgt}

\begin{orig}
Growth mixture
models\textsuperscript{\citeproc{ref-muthenshedden1999}{33}--\citeproc{ref-muthen2004}{35}}
extend LPA to longitudinal data by specifying class-specific growth
curves (or, in the trial setting, drug-response curves) with
within-class random-effects variation. A GMM is a longitudinal linear
mixed model in which the population distribution of random effects is a
mixture of multivariate normals, each corresponding to a distinct
trajectory class. Group-based trajectory
models\textsuperscript{\citeproc{ref-nagin1999}{36},\citeproc{ref-nagin2005}{37}}
are the degenerate case in which within-class random-effect variances
are set to zero, so that participants within a class share an identical
trajectory up to observation-level residual noise.

\end{orig}

\begin{bullets}

\textbf{GMM is the closest analogue of the faithful DGP in Section 3,
but Bauer and Curran's overextraction cautions apply with full force
at N-of-1 sample sizes.}

\begin{itemize}\setlength{\itemsep}{2pt}\setlength{\parskip}{0pt}
\item A biomarker-gated GMM is essentially the faithful latent-class
  DGP with drug-state-invariant class labels.
\item GMM can extract apparent class structure from single-class
  non-gaussian data at realistic sample sizes.
\item N-of-1 participant counts are typically well below those at
  which GMM class-count recovery has been validated.
\end{itemize}

\end{bullets}

\begin{rgt}
Lorem ipsum dolor sit amet, consectetur adipiscing elit, sed do eiusmod
tempor incididunt ut labore et dolore magna aliqua.

\end{rgt}

\begin{orig}
GMM is the closest psychometric analogue of the latent-class generative
form developed above: each participant belongs to a latent class with
its own drug-response trajectory, and the class label is
drug-state-invariant. A GMM that uses the biomarker as a
class-membership predictor is essentially the faithful generative model
written down in Section 3 of this paper. Bauer and
Curran\textsuperscript{\citeproc{ref-bauercurran2003}{1}} and
Bauer\textsuperscript{\citeproc{ref-bauer2007}{38}} document the
identifiability and overextraction hazards that accompany GMM estimation
at realistic sample sizes, showing in particular that GMM can extract
apparent class structure from data generated by single-class
non-gaussian processes. These cautions apply with full force to
aggregated N-of-1 trials, where participant counts are typically well
below the sample sizes at which GMM class-count recovery has been
validated.

\end{orig}

\subsection{Factor mixture models}\label{factor-mixture-models}

\begin{bullets}

\textbf{Factor mixture models combine discrete class structure with
within-class continuous factor heterogeneity, subsuming pure factor
analysis and pure LCA as limiting cases.}

\begin{itemize}\setlength{\itemsep}{2pt}\setlength{\parskip}{0pt}
\item Within-class variation is governed by a factor model;
  between-class variation by class-specific means, loadings, or
  residual covariances.
\item The pure continuous (single-class factor) and pure discrete
  (LCA) cases are special cases of the FMM framework.
\item FMMs match a wider range of biological scenarios at the cost
  of higher estimation requirements.
\end{itemize}

\end{bullets}

\begin{rgt}
Lorem ipsum dolor sit amet, consectetur adipiscing elit, sed do eiusmod
tempor incididunt ut labore et dolore magna aliqua.

\end{rgt}

\begin{orig}
Factor mixture
models\textsuperscript{\citeproc{ref-lubkemuthen2005}{2},\citeproc{ref-yung1997}{24},\citeproc{ref-arminger1999}{25}}
combine continuous latent factors with discrete latent classes, yielding
a generative structure in which within-class variation is governed by a
factor model and between-class variation is governed by class-specific
means and, optionally, class-specific factor loadings or residual
covariances. The FMM framework subsumes both the pure continuous case
(single-class factor model) and the pure discrete case (LCA with no
within-class structure) as limits.

\end{orig}

\begin{bullets}

\textbf{In FMM, the biomarker can function simultaneously as a noisy
class indicator via the gating function and a continuous within-class
modulator via factor loadings.}

\begin{itemize}\setlength{\itemsep}{2pt}\setlength{\parskip}{0pt}
\item The upper level assigns participants to classes as in LCA; the
  lower level gives each class its own random-effects structure.
\item The biological hypothesis is distinct responder classes with
  additional continuous heterogeneity within each class.
\item The hybrid representation is more biologically flexible than
  either pure LCA or pure factor analysis.
\end{itemize}

\end{bullets}

\begin{rgt}
Lorem ipsum dolor sit amet, consectetur adipiscing elit, sed do eiusmod
tempor incididunt ut labore et dolore magna aliqua.

\end{rgt}

\begin{orig}
FMMs are two-level hierarchies. At the upper level, participants are
assigned to one of a small number of classes, as in LCA. At the lower
level, each class has its own within-class random-effects structure. The
biological hypothesis the FMM encodes is that there are genuinely
distinct responder classes and that, within class, additional continuous
heterogeneity is driven by covariates or unmeasured individual
differences. For the biomarker-moderation problem, FMM provides a
principled way to represent a biomarker as both a noisy class indicator
(via the gating function) and a continuous modulator of within-class
response (via factor loadings on the response factor). The hybrid
representation matches more biological scenarios than either pure LCA or
pure factor analysis, at the cost of correspondingly higher estimation
requirements.

\end{orig}

\subsection{Regression mixtures and mixtures of
experts}\label{regression-mixtures-and-mixtures-of-experts}

\begin{bullets}

\textbf{Regression mixtures replace a single regression surface with
class-specific surfaces governed by a covariate-dependent gating
function, paralleling propensity-score stratification in causal
inference.}

\begin{itemize}\setlength{\itemsep}{2pt}\setlength{\parskip}{0pt}
\item Each class has its own regression surface; the gating model
  governs which surface applies to each participant.
\item The biomarker enters both the gating function and (via class
  membership) the within-class regression coefficient on treatment.
\item The two-stage structure parallels propensity-score
  stratification followed by within-stratum effect estimation.
\end{itemize}

\end{bullets}

\begin{rgt}
Lorem ipsum dolor sit amet, consectetur adipiscing elit, sed do eiusmod
tempor incididunt ut labore et dolore magna aliqua.

\end{rgt}

\begin{orig}
Regression
mixtures\textsuperscript{\citeproc{ref-desarbocron1988}{28}--\citeproc{ref-jedidi1997}{30}}
and the closely related mixture-of-experts
architecture\textsuperscript{\citeproc{ref-jacobs1991}{31}} allow
class-specific regression coefficients with class-membership
probabilities (the gating function) that depend on covariates. An
ordinary linear model with a biomarker-treatment interaction posits a
single regression surface in which the treatment slope varies smoothly
with the biomarker. A regression mixture instead posits two or more
distinct regression surfaces, one per class, with a separate
covariate-dependent model governing which surface applies to which
participant. The two-stage structure (classification by the gating
model, then prediction by the within-class regression) parallels the
two-stage structure of causal-inference methods that use propensity
scores to stratify before estimating within-stratum treatment effects.

\end{orig}

\begin{bullets}

\textbf{In the trial simulation idiom, regression mixtures implement a
DGP where the biomarker governs both responder probability and
within-responder effect magnitude; \texttt{flexmix} provides the R
implementation.}

\begin{itemize}\setlength{\itemsep}{2pt}\setlength{\parskip}{0pt}
\item The biomarker enters the gating function and determines
  within-responder drug effect magnitude.
\item \texttt{flexmix} implements this family with user-definable
  component models including mixed-effects specifications.
\item This family is suitable for repeated-measures DGP construction
  in aggregated N-of-1 simulation.
\end{itemize}

\end{bullets}

\begin{rgt}
Lorem ipsum dolor sit amet, consectetur adipiscing elit, sed do eiusmod
tempor incididunt ut labore et dolore magna aliqua.

\end{rgt}

\begin{orig}
In the trial-simulation idiom, regression mixtures correspond to a DGP
in which the biomarker governs both the probability of being a responder
and the magnitude of the within-responder drug effect. The
\texttt{flexmix} package in
R\textsuperscript{\citeproc{ref-leisch2004}{43},\citeproc{ref-grunleisch2008}{44}}
implements this family with user-definable component models, including
mixed-effects components suitable for repeated-measures data.

\end{orig}

\section{The biomarker-moderation
spectrum}\label{the-biomarker-moderation-spectrum}

\begin{bullets}

\textbf{Locating the covariance-only formulation within the broader
DGP spectrum clarifies which features of the carryover-induced power
loss are properties of the specification rather than the science.}

\begin{itemize}\setlength{\itemsep}{2pt}\setlength{\parskip}{0pt}
\item The spectrum is organized by which moments the biomarker
  moderates: mean, covariance, both, or higher moments.
\item Covariance-only is the Hendrickson formulation; mean-and-
  covariance and higher-moment forms are the generalizations.
\item Understanding the spectrum is prerequisite to a complete
  sensitivity analysis of carryover-induced power loss.
\end{itemize}

\end{bullets}

\begin{rgt}
Lorem ipsum dolor sit amet, consectetur adipiscing elit, sed do eiusmod
tempor incididunt ut labore et dolore magna aliqua.

\end{rgt}

\begin{orig}
The covariance-only formulation
of\textsuperscript{\citeproc{ref-hendrickson2020}{9}} specifies a
biomarker-dependent covariance structure with population means held
constant across biomarker values. Locating this formulation within the
broader spectrum of biomarker-moderated DGPs is useful because the
moment-by-moment organization of the spectrum (biomarker moderating the
mean, the covariance, both, or higher moments) clarifies which features
of the carryover-induced power loss are properties of the specification
itself rather than of the underlying scientific question.

\end{orig}

\subsection{Covariance-only
moderation}\label{covariance-only-moderation}

\begin{bullets}

\textbf{The covariance-only specification fixes the conditional mean
and encodes the interaction entirely in how strongly the biomarker
co-varies with response during on-drug versus off-drug periods.}

\begin{itemize}\setlength{\itemsep}{2pt}\setlength{\parskip}{0pt}
\item High- and low-biomarker participants share the same expected
  drug response but differ in their biomarker-response covariance.
\item The interaction signal resides entirely in the second moment
  and is eroded directly by carryover.
\item This is the core DGP in Hendrickson et al. (2020).
\end{itemize}

\end{bullets}

\begin{rgt}
Lorem ipsum dolor sit amet, consectetur adipiscing elit, sed do eiusmod
tempor incididunt ut labore et dolore magna aliqua.

\end{rgt}

\begin{orig}
The covariance-only specification holds \(E[BR_{it}]\) fixed across
biomarker strata and encodes the interaction entirely in
\(\mathrm{Cov}(B_i, BR_{it} \mid D_{bc,it})\). High- and low-biomarker
participants have the same expected drug response at any given time but
differ in the covariance between biomarker and response, specifically in
how strongly the biomarker co-varies with response during on-drug versus
off-drug periods.

\end{orig}

\begin{bullets}

\textbf{Covariance-only is an unusual restriction in the psychometric
tradition; its tractability for MVN power simulation is its main
appeal, but it lacks a clean natural generative mechanism.}

\begin{itemize}\setlength{\itemsep}{2pt}\setlength{\parskip}{0pt}
\item FMM and GMM parameterizations typically allow class-specific
  means at minimum; the means-constant-covariance-varies case tests
  measurement invariance, not primary DGP structure.
\item The encoding's appeal is tractability: the entire DGP reduces
  to a single MVN draw with a constructed covariance.
\item Its cost is the absence of a natural mechanistic
  interpretation.
\end{itemize}

\end{bullets}

\begin{rgt}
Lorem ipsum dolor sit amet, consectetur adipiscing elit, sed do eiusmod
tempor incididunt ut labore et dolore magna aliqua.

\end{rgt}

\begin{orig}
This is an unusual restriction within the psychometric mixture
tradition. Most standard FMM and GMM parameterizations permit
class-specific means at a minimum, and the means-constant-covariance-
varies case is typically encountered only as a constrained submodel used
for testing measurement-invariance
hypotheses\textsuperscript{\citeproc{ref-meredith1993}{49},\citeproc{ref-widamanreise1997}{50}}.
The corresponding biostatistical model has two latent classes with
identical expected outcomes but differing residual variance or
within-subject autocorrelation; such models are occasionally fitted as
sensitivity analyses but rarely assumed as the primary DGP. The
covariance-only encoding's appeal lies in its tractability for power
simulation under multivariate normality: the entire DGP reduces to
drawing from a single MVN with a constructed covariance. Its cost is
that it does not correspond cleanly to any natural generative mechanism.

\end{orig}

\subsection{Combined mean and covariance
moderation}\label{combined-mean-and-covariance-moderation}

\begin{bullets}

\textbf{The Arminger mean-and-covariance form is the general
psychometric parameterization, subsuming both pure mean moderation and
covariance-only as special cases.}

\begin{itemize}\setlength{\itemsep}{2pt}\setlength{\parskip}{0pt}
\item Both the expected-response vector and the residual covariance
  matrix may differ between classes and depend on covariates.
\item Mean-moderation only and covariance-only are limiting cases.
\item This is the general form against which both architectures A
  and B are constrained submodels.
\end{itemize}

\end{bullets}

\begin{rgt}
Lorem ipsum dolor sit amet, consectetur adipiscing elit, sed do eiusmod
tempor incididunt ut labore et dolore magna aliqua.

\end{rgt}

\begin{orig}
Arminger and colleagues\textsuperscript{\citeproc{ref-arminger1999}{25}}
develop finite mixtures of conditional mean- and covariance-structure
models in which both the expected-response vector and the residual
covariance matrix differ between classes and may depend on observed
covariates within class. This is the general psychometric form: it
reduces to a mean-moderation model in the limit of a single class with
covariate-dependent mean only, and to the covariance-only encoding in
the limit of class-invariant means with covariate-dependent covariance
only.

\end{orig}

\begin{bullets}

\textbf{Under combined moderation, the mean channel is robust to
carryover and the covariance channel is eroded, predicting intermediate
power loss and making this the most biologically plausible regime.}

\begin{itemize}\setlength{\itemsep}{2pt}\setlength{\parskip}{0pt}
\item Both channels carry class-membership information; the mean
  channel via $E[BR \mid B]$, the covariance channel via
  $\mathrm{Var}[BR \mid B]$.
\item Carryover attenuates the covariance channel but not the mean
  channel, yielding intermediate power loss.
\item The combined specification is plausibly the most biologically
  faithful regime for biomarker-guided psychiatric trials.
\end{itemize}

\end{bullets}

\begin{rgt}
Lorem ipsum dolor sit amet, consectetur adipiscing elit, sed do eiusmod
tempor incididunt ut labore et dolore magna aliqua.

\end{rgt}

\begin{orig}
In the combined form, responders and non-responders differ both in
typical drug response (mean moderation) and in how noisily that response
is expressed within class (covariance moderation). Both channels carry
information about class membership: the mean channel via
\(E[BR \mid B]\), the covariance channel via
\(\mathrm{Var}[BR \mid B]\). Under carryover, the mean channel is
comparatively robust because class-conditional means need not converge
off-drug; the covariance channel is eroded by construction because
off-drug covariance is attenuated. A combined-specification simulation
should therefore exhibit power loss intermediate between the pure-mean
and pure-covariance extremes; this is plausibly the most biologically
faithful regime for biomarker-guided psychiatric trials. Verbeke and
Lesaffre\textsuperscript{\citeproc{ref-verbekelesaffre1996}{40}} provide
the analogous treatment for linear mixed-effects models with
heterogeneity in the random-effects population, which is the more direct
generalization of the repeated-measures structure used in N-of-1 trial
simulation.

\end{orig}

\subsection{Location-scale moderation}\label{location-scale-moderation}

\begin{bullets}

\textbf{GAMLSS-style distributional regression allows the biomarker
to enter the mean, variance, and higher distributional moments
simultaneously via separate link functions.}

\begin{itemize}\setlength{\itemsep}{2pt}\setlength{\parskip}{0pt}
\item Standard regression models only the conditional mean; GAMLSS
  models the mean, variance, skewness, and kurtosis each as a
  separate function of covariates.
\item The familiar special case is heteroscedastic linear regression
  with covariate-dependent residual variance.
\item GAMLSS generalizes this to arbitrary distributional shape.
\end{itemize}

\end{bullets}

\begin{rgt}
Lorem ipsum dolor sit amet, consectetur adipiscing elit, sed do eiusmod
tempor incididunt ut labore et dolore magna aliqua.

\end{rgt}

\begin{orig}
Generalised additive models for location, scale, and
shape\textsuperscript{\citeproc{ref-gamlss2005}{51}} and the related
distributional regression
tradition\textsuperscript{\citeproc{ref-klein2015}{52}} allow the
biomarker to enter multiple parameters of the outcome distribution
simultaneously through separate link functions. In plain terms, a
standard regression models the conditional mean of an outcome as a
function of covariates; a GAMLSS-style distributional regression models
the conditional mean, the conditional variance, and if appropriate the
conditional skewness and kurtosis, each as a separate function of
covariates. The familiar special case is heteroscedastic linear
regression with covariate-dependent residual variance; GAMLSS
generalizes this to arbitrary distributional shape.

\end{orig}

\begin{bullets}

\textbf{GAMLSS occupies an intermediate position between the mean-moderation architecture
and B, allowing joint mean and scale moderation without discrete class
structure.}

\begin{itemize}\setlength{\itemsep}{2pt}\setlength{\parskip}{0pt}
\item A GAMLSS DGP allows a shifted mean (the mean-moderation architecture location
  effect) and altered variability (the covariance-moderation architecture analogue) jointly.
\item It is the continuous-heterogeneity analogue of FMM, combining
  mean and scale effects without requiring class structure.
\item The intermediate position is attractive for sensitivity analysis
  as it avoids committing to a class count.
\end{itemize}

\end{bullets}

\begin{rgt}
Lorem ipsum dolor sit amet, consectetur adipiscing elit, sed do eiusmod
tempor incididunt ut labore et dolore magna aliqua.

\end{rgt}

\begin{orig}
In the biomarker-trial context, a GAMLSS-based DGP would allow
high-biomarker participants to have both a shifted mean drug response (a
location effect, which is the mean-moderation architecture) and altered
response variability (a scale effect, which is analogous to the
covariance-moderation architecture's covariance moderation but operating
on marginal rather than joint moments), without requiring discrete class
structure. The GAMLSS framework is thus the continuous-heterogeneity
analogue of FMM and occupies an intermediate position between the pure
mean-moderation architecture (mean only) and the pure
covariance-moderation architecture (covariance only). This intermediate
position is attractive for sensitivity analysis because it does not
require committing to a specific number of latent classes.

\end{orig}

\subsection{Heterogeneous random-effects
models}\label{heterogeneous-random-effects-models}

\begin{bullets}

\textbf{Heterogeneous random-effects models replace the single normal
random-slope distribution with a mixture, so responder and
non-responder classes have distinct treatment-slope distributions.}

\begin{itemize}\setlength{\itemsep}{2pt}\setlength{\parskip}{0pt}
\item In the standard LMM the random treatment slope is drawn from a
  single normal; the heterogeneous extension replaces it with a
  finite mixture.
\item Responders draw from a large-slope component; non-responders
  from a near-zero-slope component.
\item The biomarker predicts which mixture component a participant
  comes from.
\end{itemize}

\end{bullets}

\begin{rgt}
Lorem ipsum dolor sit amet, consectetur adipiscing elit, sed do eiusmod
tempor incididunt ut labore et dolore magna aliqua.

\end{rgt}

\begin{orig}
Verbeke and
Lesaffre\textsuperscript{\citeproc{ref-verbekelesaffre1996}{40}}, Zhang
and Davidian\textsuperscript{\citeproc{ref-zhangdavidian2001}{41}}, and
Proust-Lima and
colleagues\textsuperscript{\citeproc{ref-proustlima2017}{42}} develop
linear mixed-effects models in which the random-effects distribution is
itself a finite mixture, allowing class-specific random-intercept and
random-slope distributions. The biostatistical setting this extends is
the standard linear mixed model with random intercept and random
treatment slope, familiar from any longitudinal trial analysis using
\texttt{nlme::lme} or \texttt{lme4::lmer}. In the standard model, the
random treatment slope is drawn from a single normal distribution with
mean equal to the average treatment effect and variance capturing
individual heterogeneity. In the heterogeneous random-effects extension,
that single distribution is replaced by a mixture: a fraction of
participants are drawn from a `responder' distribution with a large mean
treatment slope and another fraction from a `non-responder' distribution
with a near-zero mean treatment slope. The biomarker predicts which
mixture component a participant comes from.

\end{orig}

\begin{bullets}

\textbf{\texttt{lcmm} is the closest available tooling for
heterogeneous random-effects N-of-1 analyses, fitting a mixture on the
random-effects distribution and returning posterior class probabilities.}

\begin{itemize}\setlength{\itemsep}{2pt}\setlength{\parskip}{0pt}
\item The framework represents responder heterogeneity at the
  participant-specific random treatment effect -- the most clinically
  transparent representation.
\item The biomarker predicts class membership rather than entering
  the fixed-effect mean structure.
\item \texttt{lcmm} accepts standard LMM syntax and is the
  lowest-friction extension of the existing simulation framework.
\end{itemize}

\end{bullets}

\begin{rgt}
Lorem ipsum dolor sit amet, consectetur adipiscing elit, sed do eiusmod
tempor incididunt ut labore et dolore magna aliqua.

\end{rgt}

\begin{orig}
Applied to aggregated N-of-1 trial data, this framework represents
responder heterogeneity directly at the participant-specific random
treatment effect, which is arguably the most clinically transparent
representation available: responders and non-responders have distinct
random-slope distributions on the treatment indicator, and the biomarker
predicts class membership rather than entering the analysis-model mean
structure at the fixed-effect level. The \texttt{lcmm}
package\textsuperscript{\citeproc{ref-proustlima2017}{42}} implements
this family and is the closest available tooling for the longitudinal
aggregated N-of-1 use case: it accepts standard longitudinal mixed-model
syntax, fits a mixture on the random-effects distribution, and returns
posterior class probabilities for each participant. For an extension of
an existing aggregated N-of-1 simulation framework, \texttt{lcmm} is the
lowest-friction entry point.

\end{orig}

\section{Implications for analysis and
design}\label{implications-for-analysis-and-design}

Three implications follow from locating the covariance-only encoding
within the broader spectrum.

\begin{bullets}

\textbf{The covariance-only DGP is the most restrictive case of a
richer family; sensitivity analyses based on it alone risk overstating
carryover-induced power loss.}

\begin{itemize}\setlength{\itemsep}{2pt}\setlength{\parskip}{0pt}
\item Covariance-only places the entire biomarker-outcome signal in a
  channel eroded by carryover by construction.
\item Mean-structure signal survives carryover, attenuating overall
  power loss under combined or mean-only specifications.
\item A complete sensitivity analysis should span pure mean
  moderation, pure covariance moderation, and a combined
  specification.
\end{itemize}

\end{bullets}

\begin{rgt}
Lorem ipsum dolor sit amet, consectetur adipiscing elit, sed do eiusmod
tempor incididunt ut labore et dolore magna aliqua.

\end{rgt}

\begin{orig}
First, the covariance-only parameterization is the most restrictive
operational case of a much richer family. The substantial power loss
under carryover documented
in\textsuperscript{\citeproc{ref-hendrickson2020}{9}} is a specific
property of this restrictive case and need not generalize to the
mean-plus-covariance or heterogeneous-random-effects variants that are
standard in psychometric and biostatistical practice. A sensitivity
analysis based solely on the covariance-only DGP risks overstating the
carryover-induced power loss that a biomarker-guided aggregated N-of-1
trial should expect, because the covariance-only DGP places the entire
biomarker-outcome signal in a channel that carryover erodes by
construction. If a portion of the signal flows through the mean
structure, as in pure mean moderation or in the combined specification
of\textsuperscript{\citeproc{ref-arminger1999}{25}}, that component
survives carryover and overall power loss is correspondingly attenuated.
A complete sensitivity analysis should report power under at least three
points on the spectrum: pure mean moderation, pure covariance
moderation, and a combined specification.

\end{orig}

\begin{bullets}

\textbf{Identifiability concerns bound mixture-based analysis to a
sensitivity and DGP role at N-of-1 sample sizes, not a primary
analysis role.}

\begin{itemize}\setlength{\itemsep}{2pt}\setlength{\parskip}{0pt}
\item Published FMM applications involve hundreds to thousands of
  participants; N-of-1 trials at $N \in [30,150]$ are an order of
  magnitude smaller.
\item Mixture models should not be relied on as the primary analysis
  at trial-relevant $N$.
\item Mixture analysis retains value as a sensitivity check and as
  a DGP stress-test for the primary linear-mixed analysis.
\end{itemize}

\end{bullets}

\begin{rgt}
Lorem ipsum dolor sit amet, consectetur adipiscing elit, sed do eiusmod
tempor incididunt ut labore et dolore magna aliqua.

\end{rgt}

\begin{orig}
Second, the identifiability concerns documented in the GMM and FMM
literatures\textsuperscript{\citeproc{ref-bauercurran2003}{1},\citeproc{ref-lubkemuthen2005}{2},\citeproc{ref-bauer2007}{38},\citeproc{ref-nylund2007}{39}}
bound what can be asked of mixture-based analysis models at N-of-1 trial
sample sizes. Published FMM applications typically involve several
hundred to several thousand participants; biomarker- trial applications
with \(N \in [30, 150]\) are an order of magnitude smaller. Analysis
plans built around fitting \texttt{lcmm} or \texttt{flexmix} mixture
models directly to trial data are unlikely to yield reliable
class-membership inference at trial-relevant sample sizes and should not
be relied on as the primary analysis. Mixture analysis nonetheless
retains value as a sensitivity check (does class-structured residual
pattern emerge from the pre-specified analysis?) and as a DGP for power
simulation (does the primary analysis remain robust to class-structured
departures from gaussianity?).

\end{orig}

\begin{bullets}

\textbf{The heterogeneous random-slopes form is the natural DGP
extension target: it admits latent-class generative structure while
preserving \texttt{nlme::lme} analysis-model compatibility.}

\begin{itemize}\setlength{\itemsep}{2pt}\setlength{\parskip}{0pt}
\item Random-slopes parameterization extends the DGP in a direction
  internal to the standard mixed-effects framework.
\item The analysis model remains the pre-specified linear mixed-effects
  fit, preserving the audit chain.
\item Direct comparison against the covariance-only baseline is
  possible under matched effect-size calibration.
\end{itemize}

\end{bullets}

\begin{rgt}
Lorem ipsum dolor sit amet, consectetur adipiscing elit, sed do eiusmod
tempor incididunt ut labore et dolore magna aliqua.

\end{rgt}

\begin{orig}
Third, if N-of-1 simulation infrastructure pursues a richer DGP in
subsequent work, the natural target is the heterogeneous random-slopes
form rather than a full FMM. The random-slopes parameterization
preserves the analysis model's compatibility with the existing
\texttt{nlme::lme} machinery (the analysis model continues to be fitted
as a single linear mixed model with participant-specific random
treatment slopes) while admitting the biologically faithful latent-class
generative structure at the DGP level. The simulation DGP is extended in
a direction internal to the standard mixed-effects framework, and the
analysis model remains the pre-specified linear mixed-effects fit; this
preserves the audit chain of the existing framework and permits direct
comparison against the covariance-only baseline under matched
effect-size calibration.

\end{orig}

\section{Software for mixture
analyses}\label{software-for-mixture-analyses}

R packages relevant to estimating the models surveyed above, in
approximate order of fit to the longitudinal aggregated N-of-1 use case,
are as follows.

\begin{itemize}
\tightlist
\item
  \texttt{lcmm}\textsuperscript{\citeproc{ref-proustlima2017}{42}}.
  Joint latent-class mixed models for longitudinal continuous and
  ordinal outcomes; the closest available tooling for heterogeneous
  random-effects specification. Syntax is closely related to
  \texttt{nlme::lme} with a latent-class extension.
\item
  \texttt{flexmix}\textsuperscript{\citeproc{ref-leisch2004}{43},\citeproc{ref-grunleisch2008}{44}}.
  General regression- mixture framework with user-definable component
  models; suitable for regression-mixture and mixture-of-experts
  specifications. More general than \texttt{lcmm} but requires more
  configuration.
\item
  \texttt{mclust}\textsuperscript{\citeproc{ref-mclust2016}{53}}.
  Gaussian finite mixture models with a range of class-specific
  covariance parameterizations;
  supports\textsuperscript{\citeproc{ref-arminger1999}{25}}
  mean-plus-covariance specifications non- longitudinally. Useful for
  cross-sectional latent-profile analyses of baseline biomarker panels
  but does not natively handle repeated-measures structure.
\item
  \texttt{gamlss}\textsuperscript{\citeproc{ref-gamlss2005}{51}}.
  Distributional regression with covariate-dependent location, scale,
  and shape parameters.
\item
  \texttt{poLCA}\textsuperscript{\citeproc{ref-linzerlewis2011}{54}}.
  Polytomous latent class analysis for the categorical-indicator case;
  useful if baseline class indicators (e.g., symptom-checklist items)
  are incorporated into class assignment.
\item
  \texttt{OpenMx} and \texttt{lavaan} with mixture extensions.
  Structural- equation mixture modeling for FMM-style specifications;
  more general but with steeper configuration cost than \texttt{lcmm} or
  \texttt{flexmix}.
\end{itemize}

\begin{bullets}

\textbf{Mplus is the reference implementation for GMM, FMM, and
mixture-SEM; direct comparison with the identifiability literature
occasionally requires access to it.}

\begin{itemize}\setlength{\itemsep}{2pt}\setlength{\parskip}{0pt}
\item Mplus is the standard environment for the psychometric mixture-
  model simulation literature on identifiability and overextraction.
\item The Bauer-Curran, Lubke-Muthen, and Nylund studies all use
  Mplus.
\item Reproducing published benchmarks may require Mplus access
  alongside the R implementations.
\end{itemize}

\end{bullets}

\begin{rgt}
Lorem ipsum dolor sit amet, consectetur adipiscing elit, sed do eiusmod
tempor incididunt ut labore et dolore magna aliqua.

\end{rgt}

\begin{orig}
The reference implementation outside R is
Mplus\textsuperscript{\citeproc{ref-mplus}{55}}, which remains the
standard environment for GMM, FMM, and related mixture-SEM
specifications. Much of the published psychometric simulation literature
on mixture-model identifiability and overextraction
hazards\textsuperscript{\citeproc{ref-bauercurran2003}{1},\citeproc{ref-lubkemuthen2005}{2},\citeproc{ref-nylund2007}{39}}
uses Mplus, and direct comparison with that literature occasionally
requires access to it.

\end{orig}

\section{Discussion}\label{discussion}

\begin{bullets}

\textbf{The biomarker-moderation problem sits at the intersection of
biostatistical and psychometric traditions that have not previously
engaged; recognizing the encoding choice opens both sensitivity-analysis
and analysis-strategy programs.}

\begin{itemize}\setlength{\itemsep}{2pt}\setlength{\parskip}{0pt}
\item Biostatistics has efficient within-participant estimators but
  encodes moderation exclusively as a continuous interaction.
\item Psychometrics has mature finite-mixture machinery applied only
  to cross-sectional settings well above the N-of-1 range.
\item The continuous-versus-discrete choice is a substantive
  scientific question, not a parameterization convenience.
\end{itemize}

\end{bullets}

\begin{rgt}
Lorem ipsum dolor sit amet, consectetur adipiscing elit, sed do eiusmod
tempor incididunt ut labore et dolore magna aliqua.

\end{rgt}

\begin{orig}
We begin by noting that the biomarker-treatment moderation problem in
aggregated N-of-1 trials sits at the intersection of two methodological
traditions that have, until now, only weakly engaged with each other.
The biostatistical tradition has developed efficient estimators for
within-participant treatment contrasts and the meta-analytic machinery
to combine them, but has typically encoded biomarker moderation as a
continuous interaction in fixed-effects regression. The psychometric
tradition has developed mature finite-mixture and latent-class machinery
for representing discrete subpopulation structure but has applied it
primarily to cross-sectional or observational settings with sample sizes
well above the N-of-1 range. The argument of this paper is that the
choice between continuous-heterogeneity and latent-class encodings is a
substantive scientific question, not a parameterization convenience, and
that recognizing this choice opens both a sensitivity-analysis program
(the spectrum of biomarker-moderation specifications surveyed above) and
an analysis-strategy program (class-aware mixture analyses where
identifiability permits, and analysis-strategy mitigations where it does
not).

\end{orig}

\begin{bullets}

\textbf{Two limitations bound what the paper currently provides: the
five-study simulation program has not been executed, and the
prazosin-PTSD re-analysis has not been carried out.}

\begin{itemize}\setlength{\itemsep}{2pt}\setlength{\parskip}{0pt}
\item The companion simulation plan describes but does not execute
  the empirical core of the present line of work.
\item The prazosin-PTSD dataset is the most immediate real-data
  application target.
\item Re-analysis should span at least three biomarker-moderation
  encodings and a class-aware analysis where convergence permits.
\end{itemize}

\end{bullets}

\begin{rgt}
Lorem ipsum dolor sit amet, consectetur adipiscing elit, sed do eiusmod
tempor incididunt ut labore et dolore magna aliqua.

\end{rgt}

\begin{orig}
Unfortunately, two limitations bound what the paper currently provides.
First, the simulation program that would empirically resolve the
mixture-versus-mitigation efficiency question at trial-relevant sample
sizes is described in the companion simulation plan but not yet
executed. The five-study factorial outlined there (mixture-vs-MVN, the
mean-covariance-combined triptych, heterogeneous random slopes, the
three-regime breakdown grid, and the small-\(N\) identifiability and
Type I error pre-flight) constitutes the empirical core of the present
line of work and is the next deliverable in the program. Second, the
application to real biomarker-guided trial data has not been carried
out. The prazosin-PTSD analysis
of\textsuperscript{\citeproc{ref-hendrickson2020}{9}} provides the most
immediate target; the simulation program should run in parallel with
re-analysis of that dataset under at least three biomarker-moderation
encodings and a class-aware analysis where convergence permits.

\end{orig}

\begin{bullets}

\textbf{Three open questions remain: primary usability of class-aware
analyses at small $N$, analysis half-life specification under a mixture
DGP, and whether trial design should respond to the encoding question.}

\begin{itemize}\setlength{\itemsep}{2pt}\setlength{\parskip}{0pt}
\item Whether class-aware analyses are usable as primary tests or
  only as sensitivity checks requires direct simulation in the
  trial-relevant regime.
\item The continuous-decay $D_{bc,it}$ has a clean single-component
  interpretation; its meaning under class-specific carryover
  dynamics is less obvious.
\item If a latent-class DGP is biologically more plausible, designs
  minimizing carryover may gain power the standard hybrid design
  surrenders.
\end{itemize}

\end{bullets}

\begin{rgt}
Lorem ipsum dolor sit amet, consectetur adipiscing elit, sed do eiusmod
tempor incididunt ut labore et dolore magna aliqua.

\end{rgt}

\begin{orig}
Three open methodological questions remain. The first is whether
class-aware analyses are usable as primary analyses at any realistic
N-of-1 sample size, or only as sensitivity checks. The identifiability
evidence from the broader mixture-model literature suggests the latter;
only direct simulation in the trial-relevant parameter regime can settle
this. The second is how analysis half-life should be specified when the
DGP is a mixture rather than a single MVN. The continuous-decay drug
indicator \(D_{bc,it} = \exp(-\lambda t_{sd,it})\) has a clean
interpretation under a single-component generative model; its
interpretation when the underlying generative process is
class-structured and class-conditional carryover dynamics differ between
classes is less obvious and may require analysis-side mixture extension.
The third is whether the trial design itself should respond to the
moderation-encoding question: if a latent-class generative form is
biologically more plausible than a continuous one, designs that maximize
the per-participant contrast and minimize carryover may gain power that
the standard hybrid design surrenders.

\end{orig}

\begin{bullets}

\textbf{The five-study simulation program addresses the open
questions sequentially, with pre-registered ADEMP designs, uniform
analysis wrappers, and scoped infrastructure additions.}

\begin{itemize}\setlength{\itemsep}{2pt}\setlength{\parskip}{0pt}
\item Factorial designs are pre-registered through ADEMP tables
  committed before production runs.
\item Class-aware analyses are wrapped uniformly for direct
  cross-study comparability.
\item The required DGP infrastructure additions are enumerated,
  scoped, and sequenced to support both this paper and the companion
  program.
\end{itemize}

\end{bullets}

\begin{rgt}
Lorem ipsum dolor sit amet, consectetur adipiscing elit, sed do eiusmod
tempor incididunt ut labore et dolore magna aliqua.

\end{rgt}

\begin{orig}
The five-study simulation program described in the companion plan is
structured to address each of these questions sequentially. The
factorial designs are pre-registered through
ADEMP\textsuperscript{\citeproc{ref-morris2019}{56}} tables committed
before production runs; the class-aware analyses are wrapped uniformly
so that simulation results are directly comparable across studies; and
the infrastructure required to add the necessary DGPs to the existing
simulation framework is enumerated, scoped, and sequenced to support
both this paper's revisions and the companion carryover and
biomarker-interaction papers in the wider research program.

\end{orig}

\subsection{Why the analysis model is simpler than the data-generating
process}\label{why-the-analysis-model-is-simpler-than-the-data-generating-process}

\begin{bullets}

\textbf{The asymmetry between a structured DGP and a simple four-term
analysis model is intentional: a component-matched analysis is
impractical at trial-relevant $N$.}

\begin{itemize}\setlength{\itemsep}{2pt}\setlength{\parskip}{0pt}
\item The DGP decomposes into BR, PB, and TV components; the analysis
  model is \texttt{Sx \textasciitilde{} bm + t + Dbc + bm:Dbc}.
\item A component-by-component mirror would inflate the variance of
  $\beta_{bm:D}$ at $N \in [30,150]$.
\item Four reasons are given below for why the simpler analysis is
  preferred.
\end{itemize}

\end{bullets}

\begin{rgt}
Lorem ipsum dolor sit amet, consectetur adipiscing elit, sed do eiusmod
tempor incididunt ut labore et dolore magna aliqua.

\end{rgt}

\begin{orig}
A natural reaction to the simulation program described above is that the
DGP is structured (a three-component decomposition into biological
response, placebo-belief response, and time-variant natural history
under the standard pmsimstats parameterization; or a latent-class
mixture under the alternatives in this paper) while the analysis model
is a four-term linear mixed model
\texttt{Sx \textasciitilde{} bm + t + Dbc + bm:Dbc} with random
intercept and AR(1) residual correlation. A reader newly encountering
this asymmetry is reasonable to ask why the analysis does not mirror the
DGP component-by-component, and whether doing so would improve inference
about the moderation estimand. The short answer, with brief
justification, is the following.

\end{orig}

\begin{bullets}

\textbf{Component-matched analysis is impractical at trial-relevant
$N$ for four reasons: identification requirements, parameter inflation,
convergence fragility, and robustness of the moderation estimand.}

\begin{itemize}\setlength{\itemsep}{2pt}\setlength{\parskip}{0pt}
\item Many designs supply only a subset of the contrasts needed to
  identify BR, PB, and TV components separately.
\item Each Gompertz component adds 3--4 population parameters,
  inflating $\mathrm{Var}(\hat{\beta}_{bm:D})$ at $N \in [30,150]$.
\item $\beta_{bm:D}$ is robust to BR-versus-PB attribution, making
  the simpler model adequate for the validation question.
\end{itemize}

\end{bullets}

\begin{rgt}
Lorem ipsum dolor sit amet, consectetur adipiscing elit, sed do eiusmod
tempor incididunt ut labore et dolore magna aliqua.

\end{rgt}

\begin{orig}
A component-matched analysis (a non-linear mixed-effects model fitting
Gompertz BR, PB, and TV components with participant-specific random
effects) is impractical for trial-relevant \(N\) for four reasons: the
trial design must supply the contrasts that identify each component
(open-label versus blinded for PB, on-drug versus blinded-placebo for
BR, off-drug timepoints for TV), and many designs supply only some of
them; each Gompertz component costs roughly three to four parameters at
the population level and additional random-effect parameters per
participant, which at \(N \in [30, 150]\) inflates the variance of the
moderation estimand more than it reduces residual confounding; the
non-linear fitting machinery has higher convergence-failure rates and is
more sensitive to starting values than the linear-mixed analogue; and
the moderation estimand \(\beta_{bm:D}\) is robust to the BR-versus-PB
attribution for the validation question, even though it is not robust to
the larger questions of whether the moderation is class-structured (the
present paper) or covariance-versus-mean-encoded (the
\texttt{dgp-mean-moderation-vs-mvn} companion).

\end{orig}

\begin{bullets}

\textbf{Adding a phase indicator is the minimal tractable extension
that tests whether the apparent moderation shrinks under blinding.}

\begin{itemize}\setlength{\itemsep}{2pt}\setlength{\parskip}{0pt}
\item The phase-indicator extension tests whether the drug effect and
  its moderation shrink under blinding -- a PB-mediation signature.
\item For papers where the primary estimand is the marginal
  moderation, the simple four-term model is adequate.
\item The phase-indicator extension is a candidate sensitivity
  analysis for Study 1 once the primary cells are run.
\end{itemize}

\end{bullets}

\begin{rgt}
Lorem ipsum dolor sit amet, consectetur adipiscing elit, sed do eiusmod
tempor incididunt ut labore et dolore magna aliqua.

\end{rgt}

\begin{orig}
The smallest extension that retains tractability while improving the
BR-versus-PB sensitivity of the inference is a phase indicator plus its
interaction with the drug indicator, schematically
\texttt{Sx \textasciitilde{} bm + t + Dbc + phase + Dbc:phase + bm:Dbc + bm:Dbc:phase + (1 | ptID)}.
This adds three to four parameters and tests two things directly:
whether the apparent drug effect shrinks under blinding (a signature
that the open-label drug effect was partly PB-mediated) and whether the
moderation itself shrinks under blinding (a signature that the
biomarker-treatment interaction was partly biomarker-by-PB rather than
biomarker-by-BR). For papers where pharmacological-versus-placebo
disentanglement is a secondary scientific question, this extension is
the right step; for papers where the primary estimand is the marginal
moderation, the simple
\texttt{Sx \textasciitilde{} bm + t + Dbc + bm:Dbc} analysis is
adequate. The simulation program fixed in the companion plan uses the
simple form throughout; the phase-indicator extension is a candidate
sensitivity analysis to add to Study 1 (the mixture-vs-MVN factorial)
once the primary cells are run. The document at
\texttt{docs/24-component-decomposition-pedagogy.md} develops this point
in more detail with worked numerical examples.

\end{orig}

\subsection{The latent-class DGP as a generator of apparent
subadditivity}\label{the-latent-class-dgp-as-a-generator-of-apparent-subadditivity}

\begin{bullets}

\textbf{The latent-class DGP is a natural generator of apparent
subadditivity (a combined response smaller than the sum of the
separate component responses) in a one-component analysis; the
$\gamma > 0$ finding is ambiguous between direct mechanistic
interaction and emergent class structure.}

\begin{itemize}\setlength{\itemsep}{2pt}\setlength{\parskip}{0pt}
\item No direct $BR \times PB$ interaction is required in the
  class-conditional model to generate a non-zero $\hat{\gamma}$.
\item Two mechanisms produce $\gamma > 0$: direct mechanistic coupling
  and emergent covariance from class structure.
\item The companion BR-PB-TV decomposition paper defines and tests
  both mechanisms.
\end{itemize}

\end{bullets}

\begin{rgt}
Lorem ipsum dolor sit amet, consectetur adipiscing elit, sed do eiusmod
tempor incididunt ut labore et dolore magna aliqua.

\end{rgt}

\begin{orig}
A connection worth flagging explicitly: the latent-class DGP specified
in Study 1 is a natural mechanism for generating apparent subadditivity
in a downstream one-component analysis, even when no direct \(BR \times
PB\) interaction is specified in the class-conditional generative model.
The companion methods paper at
\texttt{analysis/report/06-component-decomposition/} defines
subadditivity as a non-zero coefficient \(\gamma\) on a \(BR \times PB\)
interaction term in the conditional-mean specification of the additive
decomposition, and notes that \(\gamma > 0\) in a fitted
additive-extension analysis is ambiguous between two distinct underlying
mechanisms.

\end{orig}

\begin{bullets}

\textbf{Emergent subadditivity arises when class membership correlates
with both $BR$ and $PB$ strength, producing a non-zero marginal
$E[BR \cdot PB]$ covariance term that mimics a mechanistic
interaction.}

\begin{itemize}\setlength{\itemsep}{2pt}\setlength{\parskip}{0pt}
\item Responder-class participants having higher $BR$ and $PB$ on
  average induces a positive $E[BR \cdot PB]$ covariance via the
  unobserved class label.
\item The magnitude depends on the gating function $\pi(B)$.
\item From the additive-extension perspective this is
  indistinguishable from a $\gamma$-weighted mechanistic interaction.
\end{itemize}

\end{bullets}

\begin{rgt}
Lorem ipsum dolor sit amet, consectetur adipiscing elit, sed do eiusmod
tempor incididunt ut labore et dolore magna aliqua.

\end{rgt}

\begin{orig}
The first mechanism is direct mechanistic interaction: shared neural
substrate, ceiling effects on the symptom scale, or belief-dynamics
adaptation each produce a true non-additive coupling between drug and
placebo response. The second mechanism is emergent subadditivity from
latent-class structure: if the trial population is a mixture of
responder and non-responder classes (the formulation in Section 3 of the
present paper) and class membership correlates with both \(BR\) and
\(PB\) strength (responder-class participants have both higher \(BR\)
and higher \(PB\) on average), then the marginal expected value of
\(BR \cdot
PB\) over the unobserved class label has a non-zero covariance term
whose magnitude depends on the gating function \(\pi(B)\). From the
additive-extension analysis this looks exactly like a
\(\gamma\)-weighted interaction.

\end{orig}

\begin{bullets}

\textbf{The Study 1 \texttt{pb\_class\_correlation} axis lets the
analyst distinguish emergent from mechanistic subadditivity empirically
within the simulation output.}

\begin{itemize}\setlength{\itemsep}{2pt}\setlength{\parskip}{0pt}
\item The axis crosses PB-class-independent and PB-class-correlated
  cells at matched marginal effect size.
\item Large $\hat{\gamma}$ in the PB-class-correlated cell but small
  in the PB-class-independent cell identifies the emergent mechanism.
\item The cross-paper synthesis with the BR-PB-TV decomposition paper
  directly falsifies or confirms class-aware versus mechanistic
  interpretations of apparent subadditivity.
\end{itemize}

\end{bullets}

\begin{rgt}
Lorem ipsum dolor sit amet, consectetur adipiscing elit, sed do eiusmod
tempor incididunt ut labore et dolore magna aliqua.

\end{rgt}

\begin{orig}
The Study 1 pre-registration accordingly includes a
\texttt{pb\_class\_correlation} axis that varies
\(\mathrm{Cor}(BR, PB)\) across class labels: a PB-class-independent
cell (matching the original specification) and a PB-class-correlated
cell (generating apparent subadditivity). The Study 1 simulation output
therefore lets the analyst distinguish the two mechanisms empirically:
if the additive-extension \(\hat{\gamma}\) is large in the
PB-class-correlated cell but small in the PB-class-independent cell at
matched marginal effect size, the mechanism is emergent rather than
mechanistic. The cross-paper synthesis with the BR-PB-TV decomposition
paper is therefore not merely methodological convenience but a direct
route to falsifying or confirming a class-aware versus
mechanistic-interaction interpretation of any apparent subadditivity
finding.

\end{orig}

\subsection{Progress to date}\label{progress-to-date}

The five-study production program is pre-registered under the ADEMP
framework of Morris, White, and
Crowther\textsuperscript{\citeproc{ref-morris2019}{56}} at
\texttt{analysis/scripts/latent-class-mixture-application/\ 00-ademp-pre-registration.md};
the present manuscript reports a Study 5 pilot at 240 replicates per
cell.

Phase 1 of the simulation program has reached the
infrastructure-and-pilot milestone. Three components are committed and
exercised end-to-end on the host R session.

Pre-registration. ADEMP tables for all five
studies\textsuperscript{\citeproc{ref-morris2019}{56}} are committed at
\texttt{analysis/scripts/latent-class-mixture-application/\ 00-ademp-pre-registration.md},
with cell counts (approximately 2,300 cells across the program), seeded
random-number control, an effect-size calibration sub-study, and a
deviation-log specification. Subsequent production runs are diagnosed
against these pre-registered estimands rather than against any post-hoc
wrapper choice.

\begin{bullets}

\textbf{Phase 1 has committed three infrastructure components: ADEMP
pre-registration, the \texttt{lcmm\_analysis()} wrapper, and the
heterogeneous-random-slopes DGP.}

\begin{itemize}\setlength{\itemsep}{2pt}\setlength{\parskip}{0pt}
\item ADEMP tables cover approximately 2,300 cells across the five
  studies with seeded RNG control and a deviation-log specification.
\item The class-aware analysis wrapper uses a two-stage fit anchoring
  gridsearch starts from a single-class \texttt{hlme}.
\item The heterogeneous-random-slopes DGP adds a Bernoulli class-
  membership channel with logistic gating in the standardized
  biomarker.
\end{itemize}

\end{bullets}

\begin{rgt}
Lorem ipsum dolor sit amet, consectetur adipiscing elit, sed do eiusmod
tempor incididunt ut labore et dolore magna aliqua.

\end{rgt}

\begin{orig}
Phase 1 of the simulation program has reached the
infrastructure-and-pilot milestone. Three components are committed and
exercised end-to-end on the host R session. ADEMP tables for all five
studies are committed with approximately 2,300 cells across the program,
seeded random-number control, an effect-size calibration sub-study, and
a deviation-log specification. Subsequent production runs are diagnosed
against these pre-registered estimands rather than against any post-hoc
wrapper choice.

\end{orig}

\begin{bullets}

\textbf{The \texttt{lcmm\_analysis()} wrapper extracts two moderation
signals: the primary gating-coefficient z-test and the secondary
within-class $\beta_{bm:D}$ t-test.}

\begin{itemize}\setlength{\itemsep}{2pt}\setlength{\parskip}{0pt}
\item The two-stage fit anchors gridsearch starting values from a
  single-class \texttt{hlme} before fitting \texttt{hlme(ng=2)}.
\item The gating-coefficient z-test identifies whether the biomarker
  predicts class membership.
\item The within-class $\beta_{bm:D}$ t-test is retained for
  inferential symmetry with the linear-mixed baseline.
\end{itemize}

\end{bullets}

\begin{rgt}
Lorem ipsum dolor sit amet, consectetur adipiscing elit, sed do eiusmod
tempor incididunt ut labore et dolore magna aliqua.

\end{rgt}

\begin{orig}
Class-aware analysis wrapper. An \texttt{lcmm\_analysis()} wrapper
around \texttt{lcmm::hlme} is committed at \texttt{01-lcmm-wrapper.R}
with the same input-output contract as the existing
\texttt{lme\_analysis()}. The wrapper uses a two-stage fit (a
single-class \texttt{hlme} is fitted first to anchor starting values for
a \texttt{gridsearch}-based multi-start \texttt{hlme(ng=2)} fit) and
extracts two moderation signals: the gating-coefficient z-test on the
biomarker (the primary class-aware moderation test, identifying whether
the biomarker predicts class membership) and the within-class
\(\beta_{bm:D}\) t-test (the secondary signal, retained for inferential
symmetry with the linear-mixed baseline).

\end{orig}

\begin{bullets}

\textbf{Smoke-test results confirm the heterogeneous-random-slopes DGP
produces the expected class structure and that \texttt{lcmm}
discriminates it from the single-component MVN-DGP.}

\begin{itemize}\setlength{\itemsep}{2pt}\setlength{\parskip}{0pt}
\item Single-replicate test at $N = 70$, gating slope 1.5 produces
  46\% responders and marginal $\mathrm{Cor}(B, BR) = 0.63$.
\item The linear-mixed \texttt{bm:Dbc} test detects the signal
  ($p = 0.002$) without class awareness.
\item \texttt{lcmm} returns entropy 0.86 and BIC 170 lower than the
  matched single-component fit.
\end{itemize}

\end{bullets}

\begin{rgt}
Lorem ipsum dolor sit amet, consectetur adipiscing elit, sed do eiusmod
tempor incididunt ut labore et dolore magna aliqua.

\end{rgt}

\begin{orig}
Heterogeneous-random-slopes DGP. A wrapper around
\texttt{generateData()} that adds a Bernoulli class-membership channel
with logistic gating in the standardized biomarker and class-specific
multipliers on the BR factor at on-drug timepoints is committed at
\texttt{02-heterogeneous-slopes-dgp.R}. On a single-replicate smoke test
(\(N = 70\), gating slope 1.5, responder multiplier 1.0, non-responder
multiplier 0.0) the DGP produces 46\% responders and a marginal
\(\mathrm{Cor}(B, BR) = 0.63\) that the linear-mixed \texttt{bm:Dbc}
test detects (\(p = 0.002\)) without class awareness; the \texttt{lcmm}
fit returns a posterior class-assignment entropy of 0.86 and a BIC 170
lower than the matched single-component MVN-DGP fit, indicating that the
mixture model discriminates the two DGPs at this parameter cell.

\end{orig}

\begin{bullets}

\textbf{The 100-replicate Study 5 pre-flight pilot establishes a 32:1
per-replicate cost ratio for the class-aware versus linear-mixed
analysis, with 96\% \texttt{lcmm} convergence.}

\begin{itemize}\setlength{\itemsep}{2pt}\setlength{\parskip}{0pt}
\item Per-replicate timing: 0.05 s linear-mixed baseline; 1.57 s
  class-aware analysis.
\item \texttt{lcmm} convergence rate: 96\%; four failures arose from
  gridsearch initialization.
\item Results are checkpointed at \texttt{output/03-study5-pilot.rds}.
\end{itemize}

\end{bullets}

\begin{rgt}
Lorem ipsum dolor sit amet, consectetur adipiscing elit, sed do eiusmod
tempor incididunt ut labore et dolore magna aliqua.

\end{rgt}

\begin{orig}
Study 5 pre-flight pilot. A 100-replicate pilot under the
single-component null DGP (the covariance-moderation architecture with
\(c_{bm} = 0\), \(N = 70\), hybrid open-label (OL) design) is reported
at \texttt{output/03-study5-pilot.rds}. Per-replicate timing is 0.05 s
for the linear-mixed baseline and 1.57 s for the class-aware analysis
(32:1 cost ratio, consistent with the smoke-test estimate).
\texttt{lcmm} convergence rate is 96\%; the four non-converging
replicates encountered numerical issues during the gridsearch
initialization.

\end{orig}

Type I error at \(\alpha = 0.05\) (MCSE 0.022) varies sharply across
candidate test forms, motivating the choice of the primary class-aware
moderation test:

{\def\LTcaptype{none} % do not increment counter
\begin{longtable}[]{@{}ll@{}}
\toprule\noalign{}
Test & Rejection rate \\
\midrule\noalign{}
\endhead
\bottomrule\noalign{}
\endlastfoot
lme \texttt{bm:Dbc} Wald t-test & 0.000 \\
lcmm gating-on-\texttt{bm} Wald z (unconditional) & \textbf{0.250} \\
lcmm within-class \texttt{bm:Dbc} Wald z & 0.000 \\
BIC selects ng=2 over ng=1 & 0.010 \\
LRT chi-sq(df=4), naive reference & 0.146 \\
BIC-conditional gating procedure & 0.000 \\
\end{longtable}
}

\begin{bullets}

\textbf{The unconditional Wald-on-gating inflation is a finite-sample
identifiability artifact: the conditional Wald SE underestimates true
variability because the optimizer latches onto random class boundaries.}

\begin{itemize}\setlength{\itemsep}{2pt}\setlength{\parskip}{0pt}
\item Empirical SD of the gating coefficient (63.9) is fifty times
  the median within-replicate Wald SE (0.86).
\item The gating coefficient is centered at zero with no systematic
  bias; inflation is in the SE, not the estimate.
\item This is the overextraction artifact predicted by Bauer and Curran (2003)
  and Lubke and Muthén (2005).
\end{itemize}

\end{bullets}

\begin{rgt}
Lorem ipsum dolor sit amet, consectetur adipiscing elit, sed do eiusmod
tempor incididunt ut labore et dolore magna aliqua.

\end{rgt}

\begin{orig}
The unconditional Wald-on-gating test rejects at five times the nominal
rate. Diagnosis of the rejecting replicates shows that the empirical
standard deviation of the gating coefficient across replicates (sd =
63.9) is roughly fifty times the median within-replicate Wald SE (0.86),
and the empirical sd of the z-statistic is 1.47 against a nominal 1.0.
The gating coefficient is centered at zero (median \(-0.10\), sign
balance 46/54) and shows no systematic bias; the inflation is driven by
the conditional Wald SE under-estimating the true variability of
\(\hat{\beta}_{\text{gating}}\) at trial-relevant \(N\). This is a
finite-sample identifiability artifact predicted by
the\textsuperscript{\citeproc{ref-bauercurran2003}{1}}
and\textsuperscript{\citeproc{ref-lubkemuthen2005}{2}} overextraction
literature: when the data have no class structure but \texttt{hlme} is
forced to fit \texttt{ng=2}, the optimizer latches onto random patterns
to define a class boundary, and the resulting Wald SE does not absorb
the labeling uncertainty.

\end{orig}

\begin{bullets}

\textbf{The naive likelihood-ratio test (LRT) is anti-conservative
because the true
distribution at the $ng=1$ boundary is chi-bar-squared, not
chi-squared; the bootstrapped LRT is correct but too expensive.}

\begin{itemize}\setlength{\itemsep}{2pt}\setlength{\parskip}{0pt}
\item Rejection rate 0.146 under the naive chi-squared($df=4$)
  reference.
\item The true LRT distribution at the $ng=1$ boundary is a
  chi-bar-squared mixture (McLachlan and Peel, ch.\ 6).
\item The Nylund et al. (2007) bootstrapped LRT is the standard calibration
  but is too expensive at the class-aware simulation per-fit budget.
\end{itemize}

\end{bullets}

\begin{rgt}
Lorem ipsum dolor sit amet, consectetur adipiscing elit, sed do eiusmod
tempor incididunt ut labore et dolore magna aliqua.

\end{rgt}

\begin{orig}
The naive LRT against a chi-squared(\(df=4\)) reference is similarly
anti-conservative (rejection 0.146), because the true LRT distribution
at the \texttt{ng=1} boundary is a chi-bar-squared mixture rather than a
chi-squared (\textsuperscript{\citeproc{ref-mclachlanpeel2000}{22}},
ch.~6). The bootstrapped LRT
of\textsuperscript{\citeproc{ref-nylund2007}{39}} is the standard
calibration but is computationally expensive at the per-fit budget of a
class-aware simulation.

\end{orig}

\begin{bullets}

\textbf{The BIC-conditional gating procedure controls Type I at
$N = 70$ under the null: BIC selects $ng=2$ in 1\% of replicates and
none of those also reject the gating Wald.}

\begin{itemize}\setlength{\itemsep}{2pt}\setlength{\parskip}{0pt}
\item The procedure requires $\Delta_{\text{BIC}} > 0$ (or
  $> 6$ per Raftery 1995) before testing the gating coefficient.
\item Zero conditional procedure rejections at the null gives Type I
  = 0 at the pilot's resolution.
\item Whether the procedure has adequate power under a true two-class
  DGP is the central question for the rest of Study 5.
\end{itemize}

\end{bullets}

\begin{rgt}
Lorem ipsum dolor sit amet, consectetur adipiscing elit, sed do eiusmod
tempor incididunt ut labore et dolore magna aliqua.

\end{rgt}

\begin{orig}
The pragmatic class-aware moderation test is therefore a BIC-conditional
gating procedure: fit \texttt{ng=1} and \texttt{ng=2}, prefer
\texttt{ng=2} only if BIC supports it (\(\Delta_{\text{BIC}} > 0\), or
more conservatively \(> 6\) per Raftery 1995), and only then test the
gating coefficient. At \(N = 70\) under the null DGP, BIC selects
\texttt{ng=2} in 1\% of replicates, and zero of those replicates also
reject the gating Wald, giving the conditional procedure a Type I error
of 0 at the pilot's resolution. The conditional procedure is therefore
usable as a primary moderation test at trial-relevant \(N\). Whether the
same conditional procedure has adequate power to detect a true two-class
DGP is the central question for the rest of Study 5; pilots at non-null
cells will measure this directly.

\end{orig}

\begin{bullets}

\textbf{The earlier claim that class-aware analyses are not usable as
primary tests at small $N$ requires correction: it applies to the
unconditional Wald test, not to the BIC-conditional procedure.}

\begin{itemize}\setlength{\itemsep}{2pt}\setlength{\parskip}{0pt}
\item The corrected proposition is that a single mixture-model
  coefficient's Wald test is unreliable at small $N$.
\item A model-comparison wrapper (BIC or BLRT) is required to control
  Type I error.
\item The BIC-conditional procedure is the pragmatic primary test at
  trial-relevant $N$.
\end{itemize}

\end{bullets}

\begin{rgt}
Lorem ipsum dolor sit amet, consectetur adipiscing elit, sed do eiusmod
tempor incididunt ut labore et dolore magna aliqua.

\end{rgt}

\begin{orig}
The diagnostic distinction matters for the paper's argument. The earlier
proposition that `class-aware analyses are not usable as primary at
trial-relevant \(N\)' is too strong as stated: it applies to the
unconditional Wald-on-gating test but not to the BIC-conditional
procedure. The corrected proposition is that the unconditional Wald test
on a single mixture-model coefficient is unreliable at small \(N\) and
requires a model-comparison wrapper (BIC or BLRT) to control Type I
error.

\end{orig}

\begin{bullets}

\textbf{The full five-study program is feasible on a single
workstation in approximately one week of background compute at 8--16
cores.}

\begin{itemize}\setlength{\itemsep}{2pt}\setlength{\parskip}{0pt}
\item Study 1's class-aware arm requires approximately 340 serial
  hours, or 22--44 hours on 8--16 cores.
\item Study 5 production is the largest block at approximately 60
  hours per arm on 8 cores.
\item The calibration sub-study and Study 5 pre-flight should run
  first to finalise effect-size axes before committing to Studies
  1, 2, and 3.
\end{itemize}

\end{bullets}

\begin{rgt}
Lorem ipsum dolor sit amet, consectetur adipiscing elit, sed do eiusmod
tempor incididunt ut labore et dolore magna aliqua.

\end{rgt}

\begin{orig}
Production-budget implication. Extrapolated to the pre-registered cell
counts (Study 1 has 768 cells \(\times\) 1,000 replicates = 768,000 fits
per analysis arm), the class-aware arm of Study 1 is approximately 340
hours of serial CPU time at the pilot's mean fit time. Parallelisation
across 8--16 cores reduces this to 22--44 hours per arm. Studies 2 and 3
require similar budgets; the Study 5 production run (5,000 replicates
\(\times\) three sample sizes \(\times\) three class counts \(\times\)
three DGP families) is the largest single block at approximately 60
hours per arm on 8 cores. The full program is feasible on a single
workstation over a week of background compute, with the calibration
sub-study and Study 5 pre-flight running first to finalise effect-size
axes and identifiability cuts before committing to Studies 1, 2, and 3.

\end{orig}

\begin{bullets}

\textbf{The open question raised by the pilot is whether the BIC
selection step has sufficient sensitivity under true class structure:
Pilot 4 will test this across gating slopes and class-separation
values.}

\begin{itemize}\setlength{\itemsep}{2pt}\setlength{\parskip}{0pt}
\item A conservative null-controlling procedure that rarely fires
  under true class structure provides no value.
\item Expected pattern: BIC-conditional outperforms lme at large
  class separations; underperforms at moderate separations.
\item Pilot 4 will fit the BIC-conditional procedure to the
  heterogeneous-slopes DGP from Study 3.
\end{itemize}

\end{bullets}

\begin{rgt}
Lorem ipsum dolor sit amet, consectetur adipiscing elit, sed do eiusmod
tempor incididunt ut labore et dolore magna aliqua.

\end{rgt}

\begin{orig}
Open question raised by the pilot. With Type I error of the
BIC-conditional procedure controlled at the null at \(N = 70\), the
next-most-pressing identifiability question is whether the BIC selection
step has enough sensitivity to detect a real two-class DGP at
trial-relevant \(N\). A class-aware procedure that is conservative under
the null but rarely fires under true class structure is no more useful
than a procedure that fires too often. Pilot 4 (planned for the next
iteration) will fit the same BIC-conditional procedure to the
heterogeneous-random-slopes DGP from Study 3 (with several values of the
gating slope and class-separation \(\Delta\)) and report the power
profile of the conditional procedure alongside that of the linear-mixed
baseline. We expect that the BIC-conditional procedure will outperform
the linear-mixed \texttt{bm:Dbc} test at large class separations (where
\texttt{ng=2} BIC consistently wins) but will underperform it at
moderate class separations (where BIC is undecided and the procedure
falls back to a default that does not exploit the latent-class
structure).

\end{orig}

\subsection{Pilot results (Study 5, 240
replicates)}\label{pilot-results-study-5-240-replicates}

\begin{bullets}

\textbf{The 240-replicate Study 5 pilot crossed three gating-slope
levels and three test forms at $N = 70$, with MCSE approximately
$0.030$ at working power.}

\begin{itemize}\setlength{\itemsep}{2pt}\setlength{\parskip}{0pt}
\item Wall time 558 s on 4-core \texttt{parallel::mclapply}.
\item \texttt{lcmm} convergence: 92.9\% at the null, 98.3--98.8\% at
  non-null cells.
\item MCSE approximately $0.030$ at working power near $0.5$.
\end{itemize}

\end{bullets}

\begin{rgt}
Lorem ipsum dolor sit amet, consectetur adipiscing elit, sed do eiusmod
tempor incididunt ut labore et dolore magna aliqua.

\end{rgt}

\begin{orig}
The Study 5 prototype was rerun at 240 replicates per cell on 4-core
\texttt{parallel::mclapply} with three gating-slope levels (gs
\(\in \{0, 1.0, 1.5\}\)) and three test forms (lme \texttt{bm:Dbc} Wald,
\texttt{lcmm} gating Wald unconditional, BIC-conditional gating). Wall
time was 558 s; convergence was 92.9\% under the null and 98.3--98.8\%
at non-null cells. Per-replicate output is at
\texttt{analysis/data/quick-sim/03-latent-replicates.rds}; Morris ADEMP
summary at \texttt{03-latent-summary.txt}. MCSE at the working power
near \(0.5\) is approximately \(0.030\).

\end{orig}

\begin{bullets}

\textbf{The unconditional Wald inflates Type I to $0.197$; the BIC-
conditional procedure controls Type I but is severely under-powered,
approximating a never-reject rule under heterogeneous slopes.}

\begin{itemize}\setlength{\itemsep}{2pt}\setlength{\parskip}{0pt}
\item Unconditional Wald Type I $0.197$ (95\% CI $[0.145, 0.250]$),
  decisively above the nominal 0.05; power below the lme baseline
  at both non-null cells.
\item BIC-conditional Type I $0.004$ under the null; power $0.025$
  at gs = $1.0$ and $0.004$ at gs = $1.5$.
\item Non-monotonicity at gs = $1.5$ suggests the lme captures the
  heterogeneous-slopes signal well enough that BIC continues to
  favor $ng=1$.
\end{itemize}

\end{bullets}

\begin{rgt}
Lorem ipsum dolor sit amet, consectetur adipiscing elit, sed do eiusmod
tempor incididunt ut labore et dolore magna aliqua.

\end{rgt}

\begin{orig}
The unconditional Wald-on-gating Type I error of \(0.197\) (95\%
binomial CI \([0.145, 0.250]\)) is now decisively distinguishable from
the nominal \(0.05\) (more than five MCSE of separation) and is
statistically consistent with the 100-replicate pilot's \(0.250\) point
estimate. Power at gs = \(1.0\) (\(0.304\)) and gs = \(1.5\) (\(0.360\))
sits below the lme baseline (\(0.325\) and \(0.542\) respectively), so
the unconditional Wald buys nothing in exchange for the Type I
inflation. The BIC-conditional procedure does control Type I (\(0.004\)
under the null, with upper CI near \(0.013\)), but its power is
\(0.025\) at gs = \(1.0\) and \(0.004\) at gs = \(1.5\): under the
heterogeneous-random-slopes DGP at \(N = 70\), the BIC step selects
\texttt{ng=1} in essentially every replicate and the conditional
procedure approximates a never-reject rule. The non-monotonicity (gs =
\(1.5\) below gs = \(1.0\)) suggests that as the heterogeneous-slopes
signal strengthens, a single-class lme captures it well enough that BIC
continues to favor \texttt{ng=1} and the conditional procedure does not
fire.

\end{orig}

\begin{bullets}

\textbf{The lme \texttt{bm:Dbc} test dominates both \texttt{lcmm}
tests at this $N$ and DGP regime; a canonical 1,000-replicate run is
required to map the regime where the BIC-conditional procedure becomes
competitive.}

\begin{itemize}\setlength{\itemsep}{2pt}\setlength{\parskip}{0pt}
\item lme holds nominal Type I (0 of 240) and reaches 0.542 power
  at gs = $1.5$.
\item Neither class-aware test is competitive: the unconditional
  test is anti-conservative; the BIC-conditional test is severely
  under-powered.
\item A canonical run at $N \in \{35, 70, 100\}$ with a wider
  gating-slope grid is required to locate the competitive regime,
  plausibly at large class separations.
\end{itemize}

\end{bullets}

\begin{rgt}
Lorem ipsum dolor sit amet, consectetur adipiscing elit, sed do eiusmod
tempor incididunt ut labore et dolore magna aliqua.

\end{rgt}

\begin{orig}
The lme \texttt{bm:Dbc} test holds nominal Type I (0 of 240 under the
null) and reaches \(0.542\) power at gs = \(1.5\), dominating both
\texttt{lcmm} tests at this sample size and DGP regime. We therefore
conclude, on the evidence of this pilot, that at trial-relevant \(N\)
and under the heterogeneous-random-slopes generating mechanism examined
here, neither the unconditional Wald-on-gating nor the BIC-conditional
gating test is competitive with the linear-mixed baseline: the
unconditional test is severely anti-conservative, and the
BIC-conditional test is severely under-powered. A canonical run at 1,000
replicates per cell with a wider gating-slope grid (and at
\(N \in \{35, 70, 100\}\)) is required to map the regime in which the
BIC-conditional procedure does become competitive, plausibly at large
class separations rather than under heterogeneous random slopes.

\end{orig}

\section{Conflict of interest}\label{conflict-of-interest}

The authors declare no conflicts of interest.

\section{Funding}\label{funding}

This work received no specific funding.

\section{Data availability statement}\label{data-availability-statement}

The data and code that support the findings of this study are openly
available in the pmsimstats-ng project repository. Per-replicate Monte
Carlo outputs and ADEMP-compliant cell-level summaries are versioned
under \texttt{analysis/data/}; simulation drivers are at
\texttt{analysis/scripts/}. Exact paths for this manuscript are listed
in the Reproducibility section below.

\section{Reproducibility}\label{reproducibility}

This manuscript was rendered with R 4.4.x inside the project's zzcollab
Docker container (image hash recorded in \texttt{Dockerfile}; package
set captured in \texttt{renv.lock}). Principal packages used by the
simulation pipeline are \texttt{nlme} (continuous-time linear-mixed
analysis with \texttt{corCAR1} residual correlation), \texttt{parallel}
(4-core \texttt{mclapply} for the \texttt{lcmm}-bottlenecked run),
\texttt{data.table} (the \texttt{original-extended} simulation
collection), \texttt{furrr} and \texttt{purrr} (the \texttt{tidyverse}
simulation collection), and \texttt{rmarkdown} plus \texttt{bookdown}
for the manuscript build. The class-aware analyses additionally use
\texttt{lcmm} with \texttt{gridsearch} initialization.

Random-number generator. Each simulation driver sets
\texttt{set.seed(20260507)} at the top of the replicate loop and uses
\texttt{parallel::mc.set.seed} inside \texttt{mclapply}. Per-replicate
seeds derive from the master seed by \texttt{+\ rep\_idx}.

Data and code availability. Per-replicate simulation outputs from the
240-replicate Study 5 pilot are checkpointed at
\texttt{analysis/data/quick-sim/03-latent-replicates.rds}; Morris ADEMP
cell-level summaries at
\texttt{analysis/data/quick-sim/03-latent-summary.txt}; the effect-size
calibration table at
\texttt{analysis/data/quick-sim/03-calibration-table.rds}. Driver
scripts are under
\texttt{analysis/scripts/latent-class-mixture-application/}, including
the pilot driver \texttt{03-study5-pilot.R} (executed) and the
production driver \texttt{10-study5-production.R} (committed; the
five-study production program has not yet been run). The pre-registered
ADEMP plan is at
\texttt{analysis/scripts/latent-class-mixture-application/\ 00-ademp-pre-registration.md}.
The manuscript source, paper-specific bibliography
(\texttt{references.bib}), and SIM CSL file are at
\texttt{analysis/report/03-latent-class-mixture-application/}.

\section{References}\label{references}

\protect\phantomsection\label{refs}
\begin{CSLReferences}{0}{1}
\bibitem[\citeproctext]{ref-bauercurran2003}
\CSLLeftMargin{1. }%
\CSLRightInline{Bauer DJ, Curran PJ. Distributional assumptions of
growth mixture models: Implications for overextraction of latent
trajectory classes. \emph{Psychological Methods}. 2003;8(3):338-363.}

\bibitem[\citeproctext]{ref-lubkemuthen2005}
\CSLLeftMargin{2. }%
\CSLRightInline{Lubke GH, Muthén B. Investigating population
heterogeneity with factor mixture models. \emph{Psychological Methods}.
2005;10(1):21-39.}

\bibitem[\citeproctext]{ref-guyatt1986}
\CSLLeftMargin{3. }%
\CSLRightInline{Guyatt G, Sackett D, Taylor DW, Chong J, Roberts R,
Pugsley S. Determining optimal therapy: Randomized trials in individual
patients. \emph{New England Journal of Medicine}. 1986;314(14):889-892.}

\bibitem[\citeproctext]{ref-guyatt1988}
\CSLLeftMargin{4. }%
\CSLRightInline{Guyatt G, Keller JL, Jaeschke R, Rosenthal D, Adachi JD,
Newhouse MT. The {N}-of-1 randomized controlled trial: Clinical
usefulness. Our three-year experience. \emph{Annals of Internal
Medicine}. 1988;108(2):293-299.}

\bibitem[\citeproctext]{ref-lillie2011}
\CSLLeftMargin{5. }%
\CSLRightInline{Lillie EO, Patay B, Diamant J, Issell B, Topol EJ,
Schork NJ. The n-of-1 clinical trial: The ultimate strategy for
individualizing medicine? \emph{Personalized Medicine}.
2011;8(2):161-173.}

\bibitem[\citeproctext]{ref-senn2018}
\CSLLeftMargin{6. }%
\CSLRightInline{Senn S. Statistical pitfalls of personalized medicine.
\emph{Nature}. 2018;563(7733):619-621.}

\bibitem[\citeproctext]{ref-mirza2017}
\CSLLeftMargin{7. }%
\CSLRightInline{Mirza RD, Punja S, Vohra S, Guyatt G. The history and
development of {N}-of-1 trials. \emph{Journal of the Royal Society of
Medicine}. 2017;110(8):330-340.}

\bibitem[\citeproctext]{ref-yelland2009}
\CSLLeftMargin{8. }%
\CSLRightInline{Yelland MJ, Schluter PJ. Defining worthwhile and desired
responses to treatment of chronic low back pain: A comparison of n-of-1
trials with parallel-group trials. \emph{Journal of Evaluation in
Clinical Practice}. 2009;15(3):528-533.}

\bibitem[\citeproctext]{ref-hendrickson2020}
\CSLLeftMargin{9. }%
\CSLRightInline{Hendrickson RC, Thomas RG, Schork NJ, Raskind MA.
Optimizing aggregated {N}-of-1 trial designs for predictive biomarker
validation: Statistical methods and practical considerations.
\emph{Frontiers in Digital Health}. 2020;2:13.}

\bibitem[\citeproctext]{ref-schork2015}
\CSLLeftMargin{10. }%
\CSLRightInline{Schork NJ. Personalized medicine: Time for one-person
trials. \emph{Nature}. 2015;520(7549):609-611.}

\bibitem[\citeproctext]{ref-kent2020path}
\CSLLeftMargin{11. }%
\CSLRightInline{{Kent DM, Paulus JK, Klaveren D van, et al.} The
{P}redictive {A}pproaches to {T}reatment effect {H}eterogeneity ({PATH})
{S}tatement. \emph{Annals of Internal Medicine}. 2020;172(1):35-45.}

\bibitem[\citeproctext]{ref-spearman1904}
\CSLLeftMargin{12. }%
\CSLRightInline{Spearman C. {{``General intelligence,''} objectively
determined and measured}. \emph{American Journal of Psychology}.
1904;15(2):201-292.}

\bibitem[\citeproctext]{ref-thurstone1947}
\CSLLeftMargin{13. }%
\CSLRightInline{Thurstone LL. \emph{Multiple-Factor Analysis}.
University of Chicago Press; 1947.}

\bibitem[\citeproctext]{ref-lawleymaxwell1971}
\CSLLeftMargin{14. }%
\CSLRightInline{Lawley DN, Maxwell AE. \emph{Factor Analysis as a
Statistical Method}. 2nd ed. Butterworth; 1971.}

\bibitem[\citeproctext]{ref-lazarsfeld1950}
\CSLLeftMargin{15. }%
\CSLRightInline{Lazarsfeld PF. The logical and mathematical foundation
of latent structure analysis. In: Stouffer SA, ed. \emph{Measurement and
Prediction}. Princeton University Press; 1950:362-412.}

\bibitem[\citeproctext]{ref-lazarsfeldhenry1968}
\CSLLeftMargin{16. }%
\CSLRightInline{Lazarsfeld PF, Henry NW. \emph{Latent Structure
Analysis}. Houghton Mifflin; 1968.}

\bibitem[\citeproctext]{ref-goodman1974}
\CSLLeftMargin{17. }%
\CSLRightInline{Goodman LA. Exploratory latent structure analysis using
both identifiable and unidentifiable models. \emph{Biometrika}.
1974;61(2):215-231.}

\bibitem[\citeproctext]{ref-mccutcheon1987}
\CSLLeftMargin{18. }%
\CSLRightInline{McCutcheon AL. \emph{Latent Class Analysis}. Sage;
1987.}

\bibitem[\citeproctext]{ref-vermunt2003}
\CSLLeftMargin{19. }%
\CSLRightInline{Vermunt JK. Multilevel latent class models.
\emph{Sociological Methodology}. 2003;33(1):213-239.}

\bibitem[\citeproctext]{ref-daytonmacready1988}
\CSLLeftMargin{20. }%
\CSLRightInline{Dayton CM, Macready GB. Concomitant-variable
latent-class models. \emph{Journal of the American Statistical
Association}. 1988;83(401):173-178.}

\bibitem[\citeproctext]{ref-bandeenroche1997}
\CSLLeftMargin{21. }%
\CSLRightInline{Bandeen-Roche K, Miglioretti DL, Zeger SL, Rathouz PJ.
Latent variable regression for multiple discrete outcomes. \emph{Journal
of the American Statistical Association}. 1997;92(440):1375-1386.}

\bibitem[\citeproctext]{ref-mclachlanpeel2000}
\CSLLeftMargin{22. }%
\CSLRightInline{McLachlan GJ, Peel D. \emph{Finite Mixture Models}.
Wiley; 2000.}

\bibitem[\citeproctext]{ref-bartholomew1987}
\CSLLeftMargin{23. }%
\CSLRightInline{Bartholomew DJ. \emph{Latent Variable Models and Factor
Analysis}. Griffin; 1987.}

\bibitem[\citeproctext]{ref-yung1997}
\CSLLeftMargin{24. }%
\CSLRightInline{Yung YF. Finite mixtures in confirmatory factor-analysis
models. \emph{Psychometrika}. 1997;62(3):297-330.}

\bibitem[\citeproctext]{ref-arminger1999}
\CSLLeftMargin{25. }%
\CSLRightInline{Arminger G, Stein P, Wittenberg J. Mixtures of
conditional mean- and covariance-structure models. \emph{Psychometrika}.
1999;64(4):475-494.}

\bibitem[\citeproctext]{ref-heckmansinger1984}
\CSLLeftMargin{26. }%
\CSLRightInline{Heckman J, Singer B. A method for minimizing the impact
of distributional assumptions in econometric models for duration data.
\emph{Econometrica}. 1984;52(2):271-320.}

\bibitem[\citeproctext]{ref-lindsay1995}
\CSLLeftMargin{27. }%
\CSLRightInline{Lindsay BG. \emph{Mixture Models: Theory, Geometry, and
Applications}. Vol 5. Institute of Mathematical Statistics; 1995.}

\bibitem[\citeproctext]{ref-desarbocron1988}
\CSLLeftMargin{28. }%
\CSLRightInline{DeSarbo WS, Cron WL. A maximum likelihood methodology
for clusterwise linear regression. \emph{Journal of Classification}.
1988;5(2):249-282.}

\bibitem[\citeproctext]{ref-wedeldesarbo1995}
\CSLLeftMargin{29. }%
\CSLRightInline{Wedel M, DeSarbo WS. A mixture likelihood approach for
generalized linear models. \emph{Journal of Classification}.
1995;12(1):21-55.}

\bibitem[\citeproctext]{ref-jedidi1997}
\CSLLeftMargin{30. }%
\CSLRightInline{Jedidi K, Jagpal HS, DeSarbo WS. Finite-mixture
structural equation models for response-based segmentation and
unobserved heterogeneity. \emph{Marketing Science}. 1997;16(1):39-59.}

\bibitem[\citeproctext]{ref-jacobs1991}
\CSLLeftMargin{31. }%
\CSLRightInline{Jacobs RA, Jordan MI, Nowlan SJ, Hinton GE. Adaptive
mixtures of local experts. \emph{Neural Computation}. 1991;3(1):79-87.}

\bibitem[\citeproctext]{ref-clark2013}
\CSLLeftMargin{32. }%
\CSLRightInline{{Clark SL, Muthén B, Kaprio J, et al.} Models and
strategies for factor mixture analysis: An example concerning the
structure underlying psychological disorders. \emph{Structural Equation
Modeling}. 2013;20(4):681-703.}

\bibitem[\citeproctext]{ref-muthenshedden1999}
\CSLLeftMargin{33. }%
\CSLRightInline{Muthén B, Shedden K. Finite mixture modeling with
mixture outcomes using the {EM} algorithm. \emph{Biometrics}.
1999;55(2):463-469.}

\bibitem[\citeproctext]{ref-muthen2001}
\CSLLeftMargin{34. }%
\CSLRightInline{Muthén B. Latent variable mixture modeling. In:
Marcoulides GA, Schumacker RE, eds. \emph{New Developments and
Techniques in Structural Equation Modeling}. Erlbaum; 2001:1-33.}

\bibitem[\citeproctext]{ref-muthen2004}
\CSLLeftMargin{35. }%
\CSLRightInline{Muthén B. Latent variable analysis: Growth mixture
modeling and related techniques for longitudinal data. In: Kaplan D, ed.
\emph{Handbook of Quantitative Methodology for the Social Sciences}.
Sage; 2004:345-368.}

\bibitem[\citeproctext]{ref-nagin1999}
\CSLLeftMargin{36. }%
\CSLRightInline{Nagin DS. Analyzing developmental trajectories: A
semiparametric, group-based approach. \emph{Psychological Methods}.
1999;4(2):139-157.}

\bibitem[\citeproctext]{ref-nagin2005}
\CSLLeftMargin{37. }%
\CSLRightInline{Nagin DS. \emph{Group-Based Modeling of Development}.
Harvard University Press; 2005.}

\bibitem[\citeproctext]{ref-bauer2007}
\CSLLeftMargin{38. }%
\CSLRightInline{Bauer DJ. Observations on the use of growth mixture
models in psychological research. \emph{Multivariate Behavioral
Research}. 2007;42(4):757-786.}

\bibitem[\citeproctext]{ref-nylund2007}
\CSLLeftMargin{39. }%
\CSLRightInline{Nylund KL, Asparouhov T, Muthén BO. Deciding on the
number of classes in latent class analysis and growth mixture modeling:
A {M}onte {C}arlo simulation study. \emph{Structural Equation Modeling}.
2007;14(4):535-569.}

\bibitem[\citeproctext]{ref-verbekelesaffre1996}
\CSLLeftMargin{40. }%
\CSLRightInline{Verbeke G, Lesaffre E. A linear mixed-effects model with
heterogeneity in the random-effects population. \emph{Journal of the
American Statistical Association}. 1996;91(433):217-221.}

\bibitem[\citeproctext]{ref-zhangdavidian2001}
\CSLLeftMargin{41. }%
\CSLRightInline{Zhang D, Davidian M. Linear mixed models with flexible
distributions of random effects for longitudinal data.
\emph{Biometrics}. 2001;57(3):795-802.}

\bibitem[\citeproctext]{ref-proustlima2017}
\CSLLeftMargin{42. }%
\CSLRightInline{Proust-Lima C, Philipps V, Liquet B. Estimation of
extended mixed models using latent classes and latent processes: The {R}
package lcmm. \emph{Journal of Statistical Software}. 2017;78(2):1-56.}

\bibitem[\citeproctext]{ref-leisch2004}
\CSLLeftMargin{43. }%
\CSLRightInline{Leisch F. {F}lex{M}ix: A general framework for finite
mixture models and latent class regression in {R}. \emph{Journal of
Statistical Software}. 2004;11(8):1-18.}

\bibitem[\citeproctext]{ref-grunleisch2008}
\CSLLeftMargin{44. }%
\CSLRightInline{Grün B, Leisch F. {F}lex{M}ix version 2: Finite mixtures
with concomitant variables and varying and constant parameters.
\emph{Journal of Statistical Software}. 2008;28(4):1-35.}

\bibitem[\citeproctext]{ref-zucker1997}
\CSLLeftMargin{45. }%
\CSLRightInline{Zucker DR, Schmid CH, McIntosh MW, D'Agostino RB, Selker
HP, Lau J. Combining single patient ({N}-of-1) trials to estimate
population treatment effects and to evaluate individual patient
responses to treatment. \emph{Journal of Clinical Epidemiology}.
1997;50(4):401-410.}

\bibitem[\citeproctext]{ref-zucker2010}
\CSLLeftMargin{46. }%
\CSLRightInline{Zucker DR, Ruthazer R, Schmid CH. Individual ({N}-of-1)
trials can be combined to give population comparative treatment effect
estimates: Methodologic considerations. \emph{Journal of Clinical
Epidemiology}. 2010;63(12):1312-1323.}

\bibitem[\citeproctext]{ref-gabler2011}
\CSLLeftMargin{47. }%
\CSLRightInline{Gabler NB, Duan N, Vohra S, Kravitz RL. {N}-of-1 trials
in the medical literature: A systematic review. \emph{Medical Care}.
2011;49(8):761-768.}

\bibitem[\citeproctext]{ref-clogg1995}
\CSLLeftMargin{48. }%
\CSLRightInline{Clogg CC. Latent class models. In: Arminger G, Clogg CC,
Sobel ME, eds. \emph{Handbook of Statistical Modeling for the Social and
Behavioral Sciences}. Plenum; 1995:311-359.}

\bibitem[\citeproctext]{ref-meredith1993}
\CSLLeftMargin{49. }%
\CSLRightInline{Meredith W. Measurement invariance, factor analysis and
factorial invariance. \emph{Psychometrika}. 1993;58(4):525-543.}

\bibitem[\citeproctext]{ref-widamanreise1997}
\CSLLeftMargin{50. }%
\CSLRightInline{Widaman KF, Reise SP. Exploring the measurement
invariance of psychological instruments: Applications in the substance
use domain. In: Bryant KJ, Windle M, West SG, eds. \emph{The Science of
Prevention}. American Psychological Association; 1997:281-324.}

\bibitem[\citeproctext]{ref-gamlss2005}
\CSLLeftMargin{51. }%
\CSLRightInline{Rigby RA, Stasinopoulos DM. Generalized additive models
for location, scale and shape. \emph{Journal of the Royal Statistical
Society Series C}. 2005;54(3):507-554.}

\bibitem[\citeproctext]{ref-klein2015}
\CSLLeftMargin{52. }%
\CSLRightInline{Klein N, Kneib T, Klasen S, Lang S. {B}ayesian
structured additive distributional regression for multivariate
responses. \emph{Journal of the Royal Statistical Society Series C}.
2015;64(4):569-591.}

\bibitem[\citeproctext]{ref-mclust2016}
\CSLLeftMargin{53. }%
\CSLRightInline{Scrucca L, Fop M, Murphy TB, Raftery AE. {m}clust 5:
Clustering, classification and density estimation using {G}aussian
finite mixture models. \emph{The R Journal}. 2016;8(1):289-317.}

\bibitem[\citeproctext]{ref-linzerlewis2011}
\CSLLeftMargin{54. }%
\CSLRightInline{Linzer DA, Lewis JB. {p}o{LCA}: An {R} package for
polytomous variable latent class analysis. \emph{Journal of Statistical
Software}. 2011;42(10):1-29.}

\bibitem[\citeproctext]{ref-mplus}
\CSLLeftMargin{55. }%
\CSLRightInline{Muthén LK, Muthén BO. \emph{{M}plus User's Guide}. 8th
ed. Muth{é}n; Muth{é}n; 1998-2017.}

\bibitem[\citeproctext]{ref-morris2019}
\CSLLeftMargin{56. }%
\CSLRightInline{Morris TP, White IR, Crowther MJ. Using simulation
studies to evaluate statistical methods. \emph{Statistics in Medicine}.
2019;38(11):2074-2102.}

\end{CSLReferences}

\vfill

\begin{center}\rule{0.5\linewidth}{0.5pt}\end{center}

\emph{Rendered on 2026-07-30 at 20:07 PDT.}\\
\emph{Source:
\texttt{\textasciitilde{}/prj/res/36-pmsimstats-ng/pmsimstats-ng/analysis/report/03-latent-class-mixture-application/report.Rmd}}

\end{document}
